% Template:     Informe/Reporte LaTeX
% Documento:    Archivo principal
% Versión:      4.0.0 (08/06/2017)
% Codificación: UTF-8
%
% Autor: Pablo Pizarro R.
%        Facultad de Ciencias Físicas y Matemáticas
%        Universidad de Chile
%        pablo.pizarro@ing.uchile.cl, ppizarror.com
%
% Manual template: [http://ppizarror.com/Template-Informe/]
% Licencia MIT:    [https://opensource.org/licenses/MIT/]

% CREACIÓN DEL DOCUMENTO
\documentclass[letterpaper,11pt]{article} % Articulo tamaño carta, 11pt
\usepackage[utf8]{inputenc}               % Codificación UTF-8
\def\templateversion{4.0.0}               % Versión del template

% INFORMACIÓN DEL DOCUMENTO
\def\nombredelinforme {Título del informe}
\def\temaatratar {Tema a tratar}

\def\autordeldocumento {Nombre del autor o grupo}
\def\nombredelcurso {Curso}
\def\codigodelcurso {CO-1234}

\def\nombreuniversidad {Universidad de Chile}
\def\nombrefacultad {Facultad de Ciencias Físicas y Matemáticas}
\def\departamentouniversidad {Departamento de la Universidad}
\def\imagendepartamento {departamentos/fcfm}
\def\imagendepartamentoescala {0.2}
\def\localizacionuniversidad {Santiago, Chile}

% INTEGRANTES, PROFESORES Y FECHAS
\def\tablaintegrantes {
\begin{minipage}{0.976\textwidth}
\begin{flushright}
\begin{tabular}{ll}
	Integrantes:
		& \begin{tabular}[t]{@{}l@{}}
			Integrante 1 \\
			Integrante 2
		\end{tabular} \\
	Profesores:
		& \begin{tabular}[t]{@{}l@{}}
			Profesor 1 \\
			Profesor 2
		\end{tabular} \\
	Auxiliares:
		& \begin{tabular}[t]{@{}l@{}}
			Auxiliar 1 \\
			Auxiliar 2
		\end{tabular} \\
	Ayudantes:
		& \begin{tabular}[t]{@{}l@{}}
			Ayudante 1 \\
			Ayudante 2
		\end{tabular} \\
	\multicolumn{2}{l}{Ayudante del laboratorio: Ayudante 1} \\
	& \\
	\multicolumn{2}{l}{Fecha de realización: \today} \\
	\multicolumn{2}{l}{Fecha de entrega: \today} \\
	\multicolumn{2}{l}{\localizacionuniversidad}
\end{tabular}
\end{flushright}
\end{minipage}}

% CONFIGURACIONES
% Template:     Informe/Reporte LaTeX
% Documento:    Configuraciones del template
% Versión:      4.2.7 (09/07/2017)
% Codificación: UTF-8
%
% Autor: Pablo Pizarro R.
%        Facultad de Ciencias Físicas y Matemáticas
%        Universidad de Chile
%        pablo.pizarro@ing.uchile.cl, ppizarror.com
%
% Manual template: [http://ppizarror.com/Template-Informe/]
% Licencia MIT:    [https://opensource.org/licenses/MIT/]

% CONFIGURACIONES GENERALES
\def\addemptypagetwosides {false}  % Añade pags. en blanco al imprimir a 2 caras
\def\anumsecaddtocounter {false}   % Función para insertar títulos anum aumenta contador
\def\defaultimagefolder {images/}  % Carpeta raíz de las imágenes
\def\defaultinterline {1.0}        % Interlineado por defecto [pt]
\def\defaultnewlinesize {11.0}     % Tamaño del salto de línea [pt]
\def\fontdocument {lmodern}        % Tipografía (lmodern,arial,helvet)
\def\numberedequation {true}       % Ecuaciones con \insert... numeradas
\def\pointdecimal {false}          % Decimales con punto en vez de coma
\def\romanpageuppercase {true}     % Páginas en número romano en mayúsculas
\def\showdotontitles {true}        % Punto al final de cada número de título/subtítulo
\def\tablepadding {1.0}            % Ancho de celda de las tablas

% CONFIGURACIÓN DE LAS LEYENDAS - CAPTION
\def\captionlessmarginimage {0.1}  % Margen sup/inf de fig si no hay leyenda [cm]
\def\captionlrmargin {2.0}         % Márgenes izq/der de la leyenda [cm]
\def\captionalignment {justified}  % Alineación leyenda: justified,centered,left,right
\def\captiontbmarginfigure {9.35}  % Margen sup/inf de la leyenda en figuras [pt]
\def\captiontbmargintable {7.0}    % Margen sup/inf de la leyenda en tablas [pt]
\def\captiontextbold {false}       % Etiquetas (Código,Figura,Tabla) en negrita
\def\codecaptiontop {true}         % Leyenda arriba del código fuente
\def\figurecaptiontop {false}      % Leyenda arriba de las imágenes
\def\showsectiononcaption {false}  % Muestra el número de sección en las leyendas
\def\tablecaptiontop {true}        % Leyenda arriba de las tablas

% CONFIGURACIÓN DEL ÍNDICE
\def\indexdepth {3}                % Profundidad máxima del índice
\def\indextitlemargin {7.0}        % Margen título en índice \insertindextitle [pt]
\def\showindex {true}              % Muestra el índice
\def\showindexofcode {true}        % Muestra la lista de códigos fuente
\def\showindexofcontents {true}    % Muestra la lista de contenidos
\def\showindexoffigures {true}     % Muestra la lista de figuras
\def\showindexoftables {true}      % Muestra la lista de tablas

% CONFIGURACIÓN DE LOS COLORES DEL DOCUMENTO
\def\captioncolor {black}          % Color de la etiqueta (Código,Figura,Tabla)
\def\captiontextcolor {black}      % Color de la leyenda
\def\citecolor {black}             % Color del número de las referencias o citas
\def\highlightcolor {yellow}       % Color del subrayado con \hl
\def\indextitlecolor {black}       % Color de los títulos del índice
\def\linkcolor {black}             % Color de los links del doc.
\def\maintextcolor {black}         % Color principal del texto
\def\portraittitlecolor {black}    % Color de los títulos de la portada
\def\showborderonlinks {false}     % Color de un links por un recuadro de color
\def\subsubtitlecolor {black}      % Color de los sub-subtítulos
\def\subtitlecolor {black}         % Color de los subtítulos
\def\tablelinecolor {black}        % Color de las líneas de las tablas
\def\titlecolor {black}            % Color de los títulos
\def\urlcolor {magenta}            % Color de los enlaces web (\url,\href)

% MÁRGENES DE FIGURAS
\def\marginbottomimages {-0.2}     % Margen inferior figura [cm]
\def\marginfloatimages {-13.0}     % Margen sup. fig. float \insertimageleft/right [pt]
\def\margintopimages {0.0}         % Margen superior figura [cm]

% ANEXO, CITAS, REFERENCIAS
\def\apaciterefsep {9}             % Separación entre referencias [pt] {apacite}
\def\bibtexrefsep {9}              % Separación entre referencias [pt] {bibtex}
\def\natbibrefsep {5}              % Separación entre referencias [pt] {natbib}
\def\referencenumsection {false}   % Sección de referencias numerada
\def\sectionappendixlastchar {.}   % Caracter entre el n° de la sec. de anexo y el título
\def\showappendixsecindex {false}  % Muestra el título de la sec. de anexos en el índice
\def\showappendixsectitle {false}  % Muestra el título de la sec. de anexo en el informe
\def\stylecitereferences {bibtex}  % Estilo citas y referencias (bibtex,apacite,natbib)
\def\twocolumnreferences {false}   % Referencias en dos columnas

% CONFIGURACIÓN PORTADA Y HEADERS
\def\gradecodeonportrait {false}   % Muestra el código del curso en la portada
\def\showfooter {true}             % Muestra el footer
\def\showheadertitle {true}        % Muestra título de la sección en el header

% MÁRGENES DE PÁGINA
\def\firstpagemargintop {3.8}      % Margen superior página portada [cm]
\def\pagemarginbottom {2.7}        % Margen inferior página [cm]
\def\pagemarginleft {2.54}         % Margen izquierdo página [cm]
\def\pagemarginright {2.54}        % Margen derecho página [cm]
\def\pagemargintop {3.0}           % Margen superior página [cm]

% ESTILO Y TAMAÑO DE TÍTULOS
\def\fontsizesubsubtitle {\large}  % Tamaño sub-subtítulos
\def\fontsizesubtitle {\Large}     % Tamaño subtítulos
\def\fontsizetitle {\huge}         % Tamaño títulos
\def\fontsizetitlei {\huge}        % Tamaño títulos en el índice
\def\stylesubsubtitle {\bfseries}  % Estilo sub-subtítulos
\def\stylesubtitle {\bfseries}     % Estilo subtítulos
\def\styletitle {\bfseries}        % Estilo títulos
\def\styletitlei {\bfseries}       % Estilo títulos en el índice

% OPCIONES DEL PDF COMPILADO
\def\cfgbookmarksopenlevel {1}     % Nivel de los marcadores a mostrar (1: secciones)
\def\cfgpdfbookmarkopen {true}     % Abre el panel de marcadores al abrir el pdf
\def\cfgpdfcenterwindow {true}     % Centra la ventana del lector al abrir el pdf
\def\cfgpdfdisplaydoctitle {true}  % Muestra el título del informe como título del pdf
\def\cfgpdffitwindow {false}       % Ajusta la ventana del lector al tamaño del pdf
\def\cfgpdftoolbar {true}          % Muestra la barra de herramientas del lector pdf

% NOMBRE DE OBJETOS
\def\nameabstract {Resumen}           % Nombre del resumen-abstract
\def\nameappendixsection{Anexos}      % Nombre de la sección de anexos/apéndices
\def\nameportraitpage {Portada}       % Etiqueta página de la portada
\def\namereferences {Referencias}     % Nombre de la sección de referencias
\def\nomltappendixsection {Anexo}     % Etiqueta sección en anexo/apéndices
\def\nomltcont {Índice de Contenidos} % Nombre del índice de contenidos
\def\nomltfigure {Lista de Figuras}   % Nombre del índice de la lista de figuras
\def\nomltsrc {Lista de Códigos}      % Nombre del índice de la lista de código
\def\nomlttable {Lista de Tablas}     % Nombre del índice de la lista de tablas
\def\nomltwfigure {Figura}            % Etiqueta leyenda de las figuras
\def\nomltwsrc {Código}               % Etiqueta leyenda del código fuente
\def\nomltwtable {Tabla}              % Etiqueta leyenda de las tablas


% IMPORTACIÓN DE LIBRERÍAS
% Template:     Informe/Reporte LaTeX
% Documento:    Importación de librerías
% Versión:      4.7.3 (04/02/2018)
% Codificación: UTF-8
%
% Autor: Pablo Pizarro R.
%        Facultad de Ciencias Físicas y Matemáticas
%        Universidad de Chile
%        pablo.pizarro@ing.uchile.cl, ppizarror.com
%
% Manual template: [http://latex.ppizarror.com/Template-Informe/]
% Licencia MIT:    [https://opensource.org/licenses/MIT/]

% LIBRERÍAS IMPORTANTES
\usepackage[spanish,es-nosectiondot,es-lcroman]{babel} % Idioma
\usepackage[T1]{fontenc} % Caracteres acentuados
\usepackage{\fontdocument} % Tipografía del documento
\usepackage{ifthen} % Manejo de condicionales
\ifthenelse{\equal{\showlinenumbers}{true}}{ % Muestra los números de línea
	\usepackage[switch,columnwise,running]{lineno}}{
}

% LIBRERÍAS INDEPENDIENTES
\usepackage{amsmath}       % Librerías matemáticas
\usepackage{amssymb}       % Librerías matemáticas
\usepackage{array}         % Nuevas características a las tablas
\usepackage{bigstrut}      % Líneas horizontales en tablas
\usepackage{bm}            % Caracteres en negrita en ecuaciones
\usepackage{booktabs}      % Permite manejar elementos visuales en tablas
\usepackage{caption}       % Leyendas
\usepackage{changepage}    % Condicionales para administrar páginas
\usepackage{chngcntr}      % Añade números a las leyendas
\usepackage{color}         % Colores
\usepackage{colortbl}      % Administración de color en las tablas
\usepackage{csquotes}      % Citas y comillas
\usepackage{datetime}      % Fechas
\usepackage{floatpag}      % Maneja números de páginas
\usepackage{floatrow}      % Permite adminisrar posiciones en los caption
\usepackage{framed}        % Permite creación de recuadros
\usepackage{gensymb}       % Simbología común
\usepackage{geometry}      % Dimensiones y geometría del documento
\usepackage{graphicx}      % Propiedades extra para los gráficos
\usepackage{lipsum}        % Permite crear párrafos de prueba
\usepackage{listings}      % Permite añadir código fuente
\usepackage{longtable}     % Permite utilizar tablas en varias hojas
\usepackage{lscape}        % Modo página horizontal de página
\usepackage{mathtools}     % Permite utilizar notaciones matemáticas
\usepackage{multicol}      % Múltiples columnas
\usepackage{needspace}     % Maneja los espacios en página
\usepackage{notoccite}     % Desactiva las citas en el índice
\usepackage{pdfpages}      % Permite administrar páginas en pdf
\usepackage{physics}       % Paquete de matemáticas
\usepackage{rotating}      % Permite rotación de objetos
\usepackage{sectsty}       % Cambia el estilo de los títulos
\usepackage{selinput}      % Compatibilidad con acentos
\usepackage{setspace}      % Cambia el espacio entre líneas
\usepackage{siunitx}       % Unidades del sistema internacional
\usepackage{soul}          % Permite subrayar texto
\usepackage{subfig}        % Permite agrupar imágenes
\usepackage{textcomp}      % Simbología común
\usepackage{url}           % Permite añadir enlaces
\usepackage{wasysym}       % Contiene caracteres misceláneos
\usepackage{wrapfig}       % Permite comprimir imágenes
\usepackage{xspace}        % Adminsitra espacios en párrafos y líneas

% LIBRERÍAS CON PARÁMETROS
\usepackage[makeroom]{cancel} % Cancelar términos en fórmulas
\usepackage[inline]{enumitem} % Permite enumerar ítems
\usepackage[bottom,norule,hang]{footmisc} % Estilo pie de página
\usepackage[subfigure,titles]{tocloft} % Maneja entradas en el índice
\usepackage[pdfencoding=auto,psdextra]{hyperref} % Enlaces, referencias
\usepackage[figure,table,lstlisting]{totalcount} % Contador de objetos
\usepackage[normalem]{ulem} % Permite tachar y subrayar
\usepackage[usenames,dvipsnames]{xcolor} % Paquete de colores avanzado

% LIBRERÍAS CONDICIONALES
\ifthenelse{\equal{\showdotontitles}{true}}{ % Agrega puntos a los títulos/subtítulos
	\usepackage{secdot}
	\sectiondot{subsection}
	\sectiondot{subsubsection}}{
}
\ifthenelse{\equal{\stylecitereferences}{natbib}}{ % Formato citas natbib
	\usepackage{natbib}
}{
	\ifthenelse{\equal{\stylecitereferences}{apacite}}{ % Formato citas apacite
		\usepackage{apacite}
	}{
		\ifthenelse{\equal{\stylecitereferences}{bibtex}}{ % Formato citas bibtex
		}{
		}
	}
}
\ifthenelse{\equal{\showappendixsecindex}{true}}{ % Anexos/Apéndices
	\usepackage[toc]{appendix}
}{
	\usepackage{appendix}
}

% LIBRERÍAS DEPENDIENTES
\usepackage{bookmark}      % Administración de marcadores en pdf
\usepackage{fancyhdr}      % Encabezados y pie de páginas
\usepackage{float}         % Administrador de posiciones de objetos
\usepackage{hyperxmp}      % Etiquetas opcionales para el pdf compilado
\usepackage{multirow}      % Agrega nuevas opciones a las tablas
\usepackage{titlesec}      % Administración de títulos


% IMPORTACIÓN DE FUNCIONES
% Template:     Informe/Reporte LaTeX
% Documento:    Definición de funciones
% Versión:      3.9.1 (03/06/2017)
% Codificación: UTF-8
%
% Autor: Pablo Pizarro R.
%        Facultad de Ciencias Físicas y Matemáticas
%        Universidad de Chile
%        pablo.pizarro@ing.uchile.cl, ppizarror.com
%
% Manual template: [http://ppizarror.com/Template-Informe/]
% Licencia MIT:    [https://opensource.org/licenses/MIT/]

\newcommand{\throwerror}[2]{
	% Lanza un mensaje de error
	% 	#1	Función del error
	%	#2	Mensaje
	\errmessage{Error: \noexpand#1 #2 (linea \the\inputlineno)}
	\stop
}

\newcommand{\throwwarning}[1]{
	% Lanza un mensaje de advertencia
	%	#1	Mensaje
	\errmessage{Advertencia: #1 (linea \the\inputlineno)}
}

\newcommand{\checkvardefined}[1]{
	% Comprueba si una variable está definida
	%	#1	Variable
	\ifthenelse{\isundefined{#1}}{
		\errmessage{Error: Variable \noexpand#1 no definida}
		\stop
	}{}
}

\newcommand{\emptyvarerr}[3]{
	% Lanza un mensaje de error si una variable no ha sido definida
	% 	#1	Función del error
	%	#2	Variable
	%	#3	Mensaje
	\ifx\hfuzz#2\hfuzz
		\throwerror{#1}{#3}
	\fi
}

\newcommand{\setcaptionmargincm}[1]{
	% Cambiar el margen de los caption
	% 	#1	Margen en centímetros
	\captionsetup{margin=#1cm}
}

\newcommand{\setpagemargincm}[4]{
	% Cambia márgenes de las páginas [cm]
	% 	#1	Margen izquierdo
	%	#2	Margen superior
	%	#3	Margen derecho
	%	#4	Margen inferior
	\newgeometry{left=#1cm, top=#2cm, right=#3cm, bottom=#4cm}
}

\newcommand{\changemargin}[2]{
	% Cambia los márgenes del documento
	%	#1 Margen izquierdo
	%	#2 Margen derecho
	\emptyvarerr{\changemargin}{#1}{Margen izquierdo no definido}
	\emptyvarerr{\changemargin}{#2}{Margen derecho no definido}
	\list{}{\rightmargin#2\leftmargin#1}\item[]}
\let\endchangemargin=\endlist

\newcommand{\newp}{
	% Inserta nueva línea
	\hbadness=10000 \vspace{\defaultnewlinesize pt} \par
}

\newcommand{\newpar}[1]{
	% Insertar párrafo
	% 	#1	Párrafo
	\hbadness=10000 #1 \newp
}

\newcommand{\newparnl}[1]{
	% Insertar párrafo sin nueva línea al final
	% 	#1	Párrafo
	#1 \par
}

\newcommand{\lpow}[2]{
	% Insertar sub-índice, a_b
	% 	#1	Elemento inferior (a)
	%	#2	Elemento superior (b)
	\ensuremath{{#1}_{#2}}
}

\newcommand{\pow}[2]{
	% Insertar elevado, a^b
	% 	#1	Elemento inferior (a)
	%	#2	Elemento superior (b)
	\ensuremath{{#1}^{#2}}
}

\newcommand{\fracpartial}[2]{
	% Fracción de derivadas parciales af/ax
	% 	#1	Función a derivar (f)
	%	#2	Variable a derivar (x)
	\ensuremath{\pdv{#1}{#2}}
}

\newcommand{\fracdpartial}[2]{
	% Fracción de derivadas parciales dobles a^2f/ax^2
	% 	#1	Función a derivar (f)
	%	#2	Variable a derivar (x)
	\ensuremath{\pdv[2]{#1}{#2}}
}

\newcommand{\fracnpartial}[3]{
	% Fracción de derivadas parciales en n, a^nf/ax^n
	% 	#1	Función a derivar (f)
	%	#2	Variable a derivar (x)
	%	#3	Orden (n)
	\ensuremath{\pdv[#3]{#1}{#2}}
}

\newcommand{\fracderivat}[2]{
	% Fracción de derivadas df/dx
	% 	#1	Función a derivar (f)
	%	#2	Variable a derivar (x)
	\ensuremath{\dv{#1}{#2}}
}

\newcommand{\fracdderivat}[2]{
	% Fracción de derivadas dobles d^2/dx^2
	% 	#1	Función a derivar (f)
	%	#2	Variable a derivar (x)
	\ensuremath{\dv[2]{#1}{#2}}
}

\newcommand{\fracnderivat}[3]{
	% Fracción de derivadas en n d^nf/dx^n
	% 	#1	Función a derivar (f)
	%	#2	Variable a derivar (x)
	%	#3	Orden de la derivada (n)
	\ensuremath{\dv[#3]{#1}{#2}}
}

\newcommand{\topequal}[2]{
	% Llave superior de equivalencia
	% 	#1	Elemento a igualar
	%	#2	Igualdad
	\ensuremath{\overbrace{#1}^{\mathclap{#2}}}
}

\newcommand{\underequal}[2]{
	% Llave inferior de equivalencia
	% 	#1	Elemento a igualar
	%	#2	Igualdad
	\ensuremath{\underbrace{#1}_{\mathclap{#2}}}
}

\newcommand{\topsequal}[2]{
	% Rectángulo superior de equivalencia
	% 	#1	Elemento a igualar
	%	#2	Igualdad
	\ensuremath{\overbracket{#1}^{\mathclap{#2}}}
}

\newcommand{\undersequal}[2]{
	% Rectángulo inferior de equivalencia
	% 	#1	Elemento a igualar
	%	#2	Igualdad
	\ensuremath{\underbracket{#1}_{\mathclap{#2}}}
}

\newcommand{\itemresize}[2]{
	% Redimensiona un ítem en textwidth
	% 	#1	Tamaño del nuevo objeto (En textwidth)
	%	#2	Objeto a redimensionar
	\emptyvarerr{\itemresize}{#1}{Tamano del nuevo objeto no definido}
	\emptyvarerr{\itemresize}{#2}{Objeto a redimensionar no definido}
	\resizebox{#1\textwidth}{!}{#2}
}

\newcommand{\equationresize}[2]{
	% Redimensiona una ecuación en textwidth
	% 	#1	Tamaño del nuevo objeto (En textwidth)
	%	#2	Ecuación a redimensionar
	\emptyvarerr{\equationresize}{#1}{Tamano del nuevo objeto no definido}
	\emptyvarerr{\equationresize}{#2}{Ecuacion a redimensionar no definida}
	\resizebox{#1\textwidth}{!}{$#2$}
}

\newcommand{\sectionanum}[1]{
	% Insertar un título sin número
	% 	#1	Título
	\emptyvarerr{\sectionanum}{#1}{Titulo no definido}
	\phantomsection
	\needspace{3\baselineskip}
	\section*{#1}
	\addcontentsline{toc}{section}{#1}
	\stepcounter{section}
	\ifthenelse{\equal{\showheadertitle}{true}}{
		\markboth{#1}{}
	}{}
}

\newcommand{\sectionanumnoi}[1]{
	% Insertar un título sin número sin indexar
	% 	#1	Título
	\emptyvarerr{\sectionanumnoi}{#1}{Titulo no definido}
	\phantomsection
	\needspace{3\baselineskip}
	\section*{#1}
	\stepcounter{section}
	\ifthenelse{\equal{\showheadertitle}{true}}{
		\markboth{#1}{}
	}{}
}

\newcommand{\sectionanumheadless}[1]{
	% Insertar un título sin número sin cambiar el título del header
	% 	#1	Título
	\emptyvarerr{\sectionanumnoheadless}{#1}{Titulo no definido}
	\section*{#1}
	\addcontentsline{toc}{section}{#1}
	\stepcounter{section}
}

\newcommand{\sectionanumnoiheadless}[1]{
	% Insertar un título sin número sin indexar sin cambiar el título del header
	% 	#1	Título
	\emptyvarerr{\sectionanumnoi}{#1}{Titulo no definido}
	\section*{#1}
	\stepcounter{section}
}

\newcommand{\subsectionanum}[1]{
	% Insertar un subtítulo sin número
	% 	#1	Subtítulo
	\emptyvarerr{\subsectionanum}{#1}{Subtitulo no definido}
	\subsection*{#1}
	\addcontentsline{toc}{subsection}{#1}
	\stepcounter{subsection}
}

\newcommand{\subsectionanumnoi}[1]{
	% Insertar un subtítulo sin número sin indexar
	% 	#1	Subtítulo
	\emptyvarerr{\subsectionanumnoi}{#1}{Subtitulo no definido}
	\subsection*{#1}
	\stepcounter{subsection}
}

\newcommand{\subsubsectionanum}[1]{
	% Insertar un sub-subtítulo sin número
	% 	#1	Sub-subtítulo
	\emptyvarerr{\subsubsectionanum}{#1}{Sub-subtitulo no definido}
	\subsubsection*{#1}
	\addcontentsline{toc}{subsubsection}{#1}
	\stepcounter{subsubsection}
}

\newcommand{\subsubsectionanumnoi}[1]{
	% Insertar un sub-subtítulo sin número sin indexar
	% 	#1	Sub-subtítulo
	\emptyvarerr{\subsubsectionanumnoi}{#1}{Sub-subtitulo no definido}
	\subsubsection*{#1}
	\stepcounter{subsubsection}
}

\newcommand{\insertemptypage}{
	% Crea una página vacía
	\newpage
	\setcounter{templatepagecounter}{\thepage}
	\pagenumbering{gobble}
	\null
	\thispagestyle{empty}
	\newpage
	\pagenumbering{arabic}
	\setcounter{page}{\thetemplatepagecounter}
}

\newcommand{\insertindextitle}[2]{
	% Insertar un título en un índice
	%	#1	Título
	%	#2	Margen superior en pt. (opcional), 10pt por defecto
	\emptyvarerr{\insertindextitle}{#1}{Titulo no definido}
	\ifx\hfuzz#2\hfuzz
		\addtocontents{toc}{\protect\addvspace{\indextitlemargin pt}}
	\else
		\addtocontents{toc}{\protect\addvspace{#2 pt}}
	\fi
	\addtocontents{toc}{\noindent\hyperref[swpn]{\textbf{#1}}}
}

\newcommand{\insertequation}[2][]{
	% Insertar una ecuación
	% 	#1	Label (opcional)
	%	#2	Ecuación
	\emptyvarerr{\insertequation}{#2}{Ecuacion no definida}
	\ifthenelse{\equal{\numberedequation}{true}}{
		\vspace{-0.1cm}
		\begin{equation}
			\text{#1} #2
		\end{equation}
		\vspace{-0.26cm}
	}{
		\ifx\hfuzz#1\hfuzz
		\else
			\throwwarning{Label invalido en ecuacion sin numero}
		\fi
		\insertequationanum{#2}
	}
}

\newcommand{\insertequationanum}[1]{
	% Insertar una ecuación sin número
	%	#1	Ecuación
	\emptyvarerr{\insertequationanum}{#1}{Ecuacion no definida}
	\vspace{-0.1cm}
	\begin{equation*}
		\ensuremath{#1}
	\end{equation*}
	\vspace{-0.26cm}
}

\newcommand{\insertequationcaptioned}[3][]{
	% Insertar una ecuación con leyenda
	% 	#1	Label (opcional)
	%	#2	Ecuación
	%	#3	Leyenda
	\emptyvarerr{\insertequationcaptioned}{#2}{Ecuacion no definida}
	\ifx\hfuzz#3\hfuzz
		\insertequation[#1]{#2}
	\else
		\ifthenelse{\equal{\numberedequation}{true}}{
			\vspace{-0.10cm}
			\begin{equation}
				\text{#1} #2
			\end{equation}
			\vspace{-0.65cm}
			\begin{changemargin}{\captionlrmargin cm}{\captionlrmargin cm}
				\centering \textcolor{\captiontextcolor}{#3}
				\vspace{0.05cm}
			\end{changemargin}
			\vspace{0.05cm}
		}{
			\ifx\hfuzz#1\hfuzz
			\else
				\throwwarning{Label invalido en ecuacion sin numero}
			\fi
			\insertequationcaptionedanum{#2}{#3}
		}
	\fi
}

\newcommand{\insertequationcaptionedanum}[2]{
	% Insertar una ecuación con leyenda sin número
	%	#1	Ecuación
	%	#2	Leyenda
	\emptyvarerr{\insertequationcaptionedanum}{#1}{Ecuacion no definida}
	\ifx\hfuzz#2\hfuzz
		\insertequationanum{#1}
	\else
		\vspace{-0.10cm}
		\begin{equation*}
			\ensuremath{#1}
		\end{equation*}
		\vspace{-0.65cm}
		\begin{changemargin}{\captionlrmargin cm}{\captionlrmargin cm}
			\centering \textcolor{\captiontextcolor}{#2}
			\vspace{0.05cm}
		\end{changemargin}
		\vspace{0.05cm}
	\fi
}

\newcommand{\insertgather}[1]{
	% Insertar una ecuación con el ambiente gather
	%	#1	Ecuación
	\emptyvarerr{\insertgather}{#1}{Ecuacion no definida}
	\ifthenelse{\equal{\numberedequation}{true}}{
		\vspace{-0.4cm}
		\begin{gather}
			\ensuremath{#1}
		\end{gather}
		\vspace{-0.4cm}
	}{
		\insertgatheranum{#1}
	}
}

\newcommand{\insertgatheranum}[1]{
	% Insertar una ecuación con el ambiente gather sin número
	%	#1	Ecuación
	\emptyvarerr{\insertgatheranum}{#1}{Ecuacion no definida}
	\vspace{-0.4cm}
	\begin{gather*}
		\ensuremath{#1}
	\end{gather*}
	\vspace{-0.4cm}
}

\newcommand{\insertgathercaptioned}[2]{
	% Insertar una ecuación (gather) con leyenda
	%	#1	Ecuación
	%	#2	Leyenda
	\emptyvarerr{\insertgathercaptioned}{#1}{Ecuacion no definida}
	\ifx\hfuzz#2\hfuzz
		\insertgather{#1}
	\else
		\ifthenelse{\equal{\numberedequation}{true}}{
			\vspace{-0.45cm}
			\begin{gather}
				\ensuremath{#1}
			\end{gather}
			\vspace{-0.77cm}
			\begin{changemargin}{\captionlrmargin cm}{\captionlrmargin cm}
				\centering \textcolor{\captiontextcolor}{#2}
				\vspace{0.05cm}
			\end{changemargin}
			\vspace{0cm}
		}{
			\insertgathercaptionedanum{#1}{#2}
		}
	\fi
}

\newcommand{\insertgathercaptionedanum}[2]{
	% Insertar una ecuación (gather) con leyenda sin número
	%	#1	Ecuación
	%	#2	Leyenda
	\emptyvarerr{\insertgathercaptionedanum}{#1}{Ecuacion no definida}
	\ifx\hfuzz#2\hfuzz
		\insertgatheranum{#1}
	\else
		\vspace{-0.45cm}
		\begin{gather*}
			\ensuremath{#1}
		\end{gather*}
		\vspace{-0.77cm}
		\begin{changemargin}{\captionlrmargin cm}{\captionlrmargin cm}
			\centering \textcolor{\captiontextcolor}{#2}
			\vspace{0.05cm}
		\end{changemargin}
		\vspace{0cm}
	\fi
}

\newcommand{\insertgathered}[2][]{
	% Insertar una ecuación con el ambiente gathered
	% 	#1	Label (opcional)
	%	#2	Ecuación
	\emptyvarerr{\insertgathered}{#2}{Ecuacion no definida}
	\ifthenelse{\equal{\numberedequation}{true}}{
		\vspace{-0.1cm}
		\begin{equation}
			\begin{gathered}
				\text{#1} \ensuremath{#2}
			\end{gathered}
		\end{equation}
		\vspace{-0.05cm}
	}{
		\ifx\hfuzz#1\hfuzz
		\else
			\throwwarning{Label invalido en ecuacion (gathered) sin numero}
		\fi
		\insertgatheredanum{#2}
	}
}

\newcommand{\insertgatheredanum}[1]{
	% Insertar una ecuación con el ambiente gathered sin número
	%	#1	Ecuación
	\emptyvarerr{\insertgatheredanum}{#1}{Ecuacion no definida}
	\vspace{-0.4cm}
	\begin{gather*}
		\ensuremath{#1}
	\end{gather*}
	\vspace{-0.4cm}
}

\newcommand{\insertgatheredcaptioned}[3][]{
	% Insertar una ecuación (gathered) con leyenda
	% 	#1	Label (opcional)
	%	#2	Ecuación
	%	#3	Leyenda
	\emptyvarerr{\insertgatheredcaptioned}{#2}{Ecuacion no definida}
	\ifx\hfuzz#3\hfuzz
		\insertgathered[#1]{#2}
	\else
		\ifthenelse{\equal{\numberedequation}{true}}{
			\vspace{0cm}
			\begin{equation}
				\begin{gathered}
					\text{#1} \ensuremath{#2}
				\end{gathered}
			\end{equation}
			\vspace{-0.65cm}
			\begin{changemargin}{\captionlrmargin cm}{\captionlrmargin cm}
				\centering \textcolor{\captiontextcolor}{#3}
				\vspace{0.05cm}
			\end{changemargin}
			\vspace{0cm}
		}{
			\ifx\hfuzz#1\hfuzz
			\else
				\throwwarning{Label invalido en ecuacion (gathered) sin numero}
			\fi
			\insertgatheredcaptionedanum{#2}{#3}
		}
	\fi
}

\newcommand{\insertgatheredcaptionedanum}[2]{
	% Insertar una ecuación (gathered) con leyenda sin número
	%	#1	Ecuación
	%	#2	Leyenda
	\emptyvarerr{\insertgatheredcaptionedanum}{#1}{Ecuacion no definida}
	\ifx\hfuzz#2\hfuzz
		\insertgatheredanum{#1}
	\else
		\vspace{-0.45cm}
		\begin{gather*}
			\ensuremath{#1}
		\end{gather*}
		\vspace{-0.7cm}
		\begin{changemargin}{\captionlrmargin cm}{\captionlrmargin cm}
			\centering \textcolor{\captiontextcolor}{#2}
			\vspace{0.05cm}
		\end{changemargin}
		\vspace{0cm}
	\fi
}

\newcommand{\insertalign}[1]{
	% Insertar una ecuación con el ambiente align
	%	#1	Ecuación
	\emptyvarerr{\insertalign}{#1}{Ecuacion no definida}
	\ifthenelse{\equal{\numberedequation}{true}}{
		\vspace{-0.45cm}
		\begin{align}
			\ensuremath{#1}
		\end{align}
		\vspace{-0.4cm}
	}{
		\insertalignanum{#1}
	}
}

\newcommand{\insertalignanum}[1]{
	% Insertar una ecuación con el ambiente align sin número
	%	#1	Ecuación
	\emptyvarerr{\insertalignanum}{#1}{Ecuacion no definida}
	\vspace{-0.45cm}
	\begin{align*}
		\ensuremath{#1}
	\end{align*}
	\vspace{-0.4cm}
}

\newcommand{\insertaligncaptioned}[2]{
	% Insertar una ecuación (align) con leyenda
	%	#1	Ecuación
	%	#2	Leyenda
	\emptyvarerr{\insertaligncaptioned}{#1}{Ecuacion no definida}
	\ifx\hfuzz#2\hfuzz
		\insertalign{#1}
	\else
		\ifthenelse{\equal{\numberedequation}{true}}{
			\vspace{-0.45cm}
			\begin{align}
				\ensuremath{#1}
			\end{align}
			\vspace{-0.77cm}
			\begin{changemargin}{\captionlrmargin cm}{\captionlrmargin cm}
				\centering \textcolor{\captiontextcolor}{#2}
				\vspace{0.05cm}
			\end{changemargin}
			\vspace{0cm}
		}{
			\insertaligncaptionedanum{#1}{#2}
		}
	\fi
}

\newcommand{\insertaligncaptionedanum}[2]{
	% Insertar una ecuación (align) con leyenda sin número
	%	#1	Ecuación
	%	#2	Leyenda
	\emptyvarerr{\insertaligncaptioned}{#1}{Ecuacion no definida}
	\ifx\hfuzz#2\hfuzz
		\insertalignanum{#1}
	\else
		\vspace{-0.45cm}
		\begin{align*}
			\ensuremath{#1}
		\end{align*}
		\vspace{-0.77cm}
		\begin{changemargin}{\captionlrmargin cm}{\captionlrmargin cm}
			\centering \textcolor{\captiontextcolor}{#2}
			\vspace{0.05cm}
		\end{changemargin}
		\vspace{0cm}
	\fi
}

\newcommand{\insertaligned}[2][]{
	% Insertar una ecuación con el ambiente aligned
	% 	#1	Label (opcional)
	%	#2	Ecuación
	\emptyvarerr{\insertaligned}{#2}{Ecuacion no definida}
	\ifthenelse{\equal{\numberedequation}{true}}{
		\vspace{-0.1cm}
		\begin{equation}
			\begin{aligned}
				\text{#1} \ensuremath{#2}
			\end{aligned}
		\end{equation}
		\vspace{-0.05cm}
	}{
		\ifx\hfuzz#1\hfuzz
		\else
			\throwwarning{Label invalido en ecuacion (aligned) sin numero}
		\fi
		\insertalignedanum{#2}
	}
}

\newcommand{\insertalignedanum}[1]{
	% Insertar una ecuación con el ambiente aligned sin número
	%	#1	Ecuación
	\emptyvarerr{\insertalignedanum}{#1}{Ecuacion no definida}
	\vspace{-0.45cm}
	\begin{align*}
		\ensuremath{#1}
	\end{align*}
	\vspace{-0.4cm}
}

\newcommand{\insertalignedcaptioned}[3][]{
	% Insertar una ecuación (aligned) con leyenda
	% 	#1	Label (opcional)
	%	#2	Ecuación
	%	#3	Leyenda
	\emptyvarerr{\insertalignedcaptioned}{#2}{Ecuacion no definida}
	\ifx\hfuzz#3\hfuzz
		\insertaligned[#1]{#2}
	\else
	\ifthenelse{\equal{\numberedequation}{true}}{
		\vspace{0cm}
		\begin{equation}
			\begin{aligned}
				\text{#1} \ensuremath{#2}
			\end{aligned}
		\end{equation}
		\vspace{-0.65cm}
		\begin{changemargin}{\captionlrmargin cm}{\captionlrmargin cm}
			\centering \textcolor{\captiontextcolor}{#3}
			\vspace{0.05cm}
		\end{changemargin}
		\vspace{0cm}
	}{
		\ifx\hfuzz#1\hfuzz
		\else
			\throwwarning{Label invalido en ecuacion (aligned) sin numero}
		\fi
		\insertalignedcaptionedanum{#2}{#3}
	}
	\fi
}

\newcommand{\insertalignedcaptionedanum}[2]{
	% Insertar una ecuación (aligned) con leyenda sin número
	%	#1	Ecuación
	%	#2	Leyenda
	\emptyvarerr{\insertalignedcaptioned}{#1}{Ecuacion no definida}
	\ifx\hfuzz#2\hfuzz
		\insertalignedanum{#1}
	\else
		\vspace{0cm}
		\begin{equation}
			\begin{aligned}
				\ensuremath{#1}
			\end{aligned}
		\end{equation}
		\vspace{-0.65cm}
		\begin{changemargin}{\captionlrmargin cm}{\captionlrmargin cm}
			\centering \textcolor{\captiontextcolor}{#2}
			\vspace{0.05cm}
		\end{changemargin}
		\vspace{0cm}
	\fi
}

\newcommand{\insertimage}[4][]{
	% Insertar una imagen
	% 	#1	Label (opcional)
	%	#2	Dirección de la imagen
	%	#3	Parámetros de la imagen
	%	#4	Leyenda de la imagen
	\emptyvarerr{\insertimage}{#2}{Direccion de la imagen no definida}
	\emptyvarerr{\insertimage}{#3}{Parametros de la imagen no definidos}
	\vspace{\margintopimages cm}
	\begin{figure}[H]
		\centering
		\includegraphics[#3]{\defaultimagefolder#2}
		\ifx\hfuzz#4\hfuzz
			\vspace{\captionlessmarginimage cm}
		\else
			\hspace{0cm} \caption{#4 #1}
		\fi
	\end{figure}
	\vspace{\marginbottomimages cm}
}

\newcommand{\insertimageboxed}[4][]{
	% Insertar una imagen con recuadro
	% 	#1	Label (opcional)
	%	#2	Dirección de la imagen
	%	#3	Parámetros de la imagen
	%	#4	Leyenda de la imagen
	\emptyvarerr{\insertimageboxed}{#2}{Direccion de la imagen no definida}
	\emptyvarerr{\insertimageboxed}{#3}{Parametros de la imagen no definidos}
	\vspace{\margintopimages cm}
	\begin{figure}[H]
		\centering
		\fbox{\includegraphics[#3]{\defaultimagefolder#2}}
		\ifx\hfuzz#4\hfuzz
			\vspace{\captionlessmarginimage cm}
		\else
			\hspace{0cm} \caption{#4 #1}
		\fi
	\end{figure}
	\vspace{\marginbottomimages cm}
}

\newcommand{\insertdoubleimage}[8][]{
	% Insertar una imagen doble
	% 	#1	Label (opcional)
	%	#2	Dirección de la imagen 1
	%	#3	Parámetros de la imagen 1
	%	#4	Leyenda de la imagen 1
	%	#5	Dirección de la imagen 2
	%	#6	Parámetros de la imagen 2
	%	#7	Leyenda de la imagen 2
	%	#8	Leyenda de la imagen doble
	\emptyvarerr{\insertdoubleimage}{#2}{Direccion de la imagen 1 no definida}
	\emptyvarerr{\insertdoubleimage}{#3}{Parametros de la imagen 1 no definidos}
	\emptyvarerr{\insertdoubleimage}{#5}{Direccion de la imagen 2 no definida}
	\emptyvarerr{\insertdoubleimage}{#6}{Parametros de la imagen 2 no definidos}
	\vspace{\margintopimages cm}
	\captionsetup{margin=0.45cm}
	\begin{figure}[H] \centering
		\subfloat[#4]{
			\includegraphics[#3]{\defaultimagefolder#2}}
		\hspace{0.2cm}
		\subfloat[#7]{
			\includegraphics[#6]{\defaultimagefolder#5}}
		\setcaptionmargincm{\captionlrmargin}
		\ifx\hfuzz#8\hfuzz
			\vspace{\captionlessmarginimage cm}
		\else
			\caption{#8 #1}
		\fi
	\end{figure}
	\setcaptionmargincm{\captionlrmargin}
	\vspace{\marginbottomimages cm}
}

\newcommand{\insertdoubleeqimage}[7][]{
	% Insertar una imagen doble, igual propiedades
	% 	#1	Label (opcional)
	%	#2	Dirección de la imagen 1
	%	#3	Leyenda de la imagen 1
	%	#4	Dirección de la imagen 2
	%	#5	Leyenda de la imagen 2
	%	#6	Propiedades de las imágenes
	%	#7 	Leyenda de la imagen doble
	\insertdoubleimage[#1]{#2}{#6}{#3}{#4}{#6}{#5}{#7}
}

\newcommand{\inserttripleimage}[8][]{
	% Insertar una imagen triple
	% 	#1	Label (opcional)
	%	#2	Dirección de la imagen 1
	%	#3	Parámetros de la imagen 1
	%	#4	Dirección de la imagen 2
	%	#5	Parámetros de la imagen 2
	%	#6	Dirección de la imagen 3
	%	#7	Parámetros de la imagen 3
	%	#8	Leyenda de la imagen triple
	\emptyvarerr{\inserttripleimage}{#2}{Direccion de la imagen 1 no definida}
	\emptyvarerr{\inserttripleimage}{#3}{Parametros de la imagen 1 no definidos}
	\emptyvarerr{\inserttripleimage}{#4}{Direccion de la imagen 2 no definida}
	\emptyvarerr{\inserttripleimage}{#5}{Parametros de la imagen 2 no definidos}
	\emptyvarerr{\inserttripleimage}{#6}{Direccion de la imagen 3 no definida}
	\emptyvarerr{\inserttripleimage}{#7}{Parametros de la imagen 3 no definidos}
	\vspace{\margintopimages cm}
	\captionsetup{margin=0.45cm}
	\begin{figure}[H] \centering
		\subfloat[]{
			\includegraphics[#3]{\defaultimagefolder#2}}
		\hspace{0.1cm}
		\subfloat[]{
			\includegraphics[#5]{\defaultimagefolder#4}}
		\hspace{0.1cm}
		\subfloat[]{
			\includegraphics[#7]{\defaultimagefolder#6}}
		\setcaptionmargincm{\captionlrmargin}
		\ifx\hfuzz#8\hfuzz
			\vspace{\captionlessmarginimage cm}
		\else
			\caption{#8 #1}
		\fi
	\end{figure}
	\setcaptionmargincm{\captionlrmargin}
	\vspace{\marginbottomimages cm}
}

\newcommand{\inserttripleeqimage}[6][]{
	% Insertar una imagen triple, igual propiedades
	% 	#1	Label (opcional)
	%	#2	Dirección de la imagen 1
	%	#3	Dirección de la imagen 2
	%	#4	Dirección de la imagen 3
	%	#5	Propiedades de las imágenes
	%	#6	Leyenda de la imagen triple
	\inserttripleimage[#1]{#2}{#5}{#3}{#5}{#4}{#5}{#6}
}

\newcommand{\insertquadimage}[7][]{
	% Insertar una imagen cuádruple, igual propiedades
	% 	#1	Label (opcional)
	%	#2	Dirección de la imagen 1
	%	#3	Dirección de la imagen 2
	%	#4	Dirección de la imagen 3
	%	#5	Dirección de la imagen 4
	%	#6	Propiedades de las imágenes
	%	#7	Leyenda de la imagen cuádruple
	\emptyvarerr{\insertquadimage}{#2}{Direccion de la imagen 1 no definida}
	\emptyvarerr{\insertquadimage}{#3}{Direccion de la imagen 2 no definida}
	\emptyvarerr{\insertquadimage}{#4}{Direccion de la imagen 3 no definida}
	\emptyvarerr{\insertquadimage}{#5}{Direccion de la imagen 4 no definida}
	\emptyvarerr{\insertquadimage}{#6}{Propiedades de las imagenes no definidos}
	\vspace{\margintopimages cm}
	\captionsetup{margin=0.45cm}
	\begin{figure}[H] \centering
		\subfloat[]{
			\includegraphics[#6]{\defaultimagefolder#2}}
		\hspace{0.1cm}
		\subfloat[]{
			\includegraphics[#6]{\defaultimagefolder#3}}
		\hspace{0.1cm}
		\subfloat[]{
			\includegraphics[#6]{\defaultimagefolder#4}}
		\hspace{0.1cm}
		\subfloat[]{
			\includegraphics[#6]{\defaultimagefolder#5}}
		\setcaptionmargincm{\captionlrmargin}
		\ifx\hfuzz#7\hfuzz
			\vspace{\captionlessmarginimage cm}
		\else
			\caption{#7 #1}
		\fi
	\end{figure}
	\setcaptionmargincm{\captionlrmargin}
	\vspace{\marginbottomimages cm}
}

\newcommand{\insertpentaimage}[8][]{
	% Insertar una imagen quíntuple, igual propiedades
	% 	#1	Label (opcional)
	%	#2	Dirección de la imagen 1
	%	#3	Dirección de la imagen 2
	%	#4	Dirección de la imagen 3
	%	#5	Dirección de la imagen 4
	%	#6	Dirección de la imagen 5
	%	#7	Propiedades de las imágenes
	%	#8	Leyenda de la imagen quíntuple
	\emptyvarerr{\insertpentaimage}{#2}{Direccion de la imagen 1 no definida}
	\emptyvarerr{\insertpentaimage}{#3}{Direccion de la imagen 2 no definida}
	\emptyvarerr{\insertpentaimage}{#4}{Direccion de la imagen 3 no definida}
	\emptyvarerr{\insertpentaimage}{#5}{Direccion de la imagen 4 no definida}
	\emptyvarerr{\insertpentaimage}{#6}{Direccion de la imagen 5 no definida}
	\emptyvarerr{\insertpentaimage}{#7}{Propiedades de las imagenes no definidas}
	\vspace{\margintopimages cm}
	\captionsetup{margin=0.45cm}
	\begin{figure}[H] \centering
		\subfloat[]{
			\includegraphics[#7]{\defaultimagefolder#2}}
		\hspace{0.1cm}
		\subfloat[]{
			\includegraphics[#7]{\defaultimagefolder#3}}
		\hspace{0.1cm}
		\subfloat[]{
			\includegraphics[#7]{\defaultimagefolder#4}}
		\hspace{0.1cm}
		\subfloat[]{
			\includegraphics[#7]{\defaultimagefolder#5}}
		\hspace{0.1cm}
		\subfloat[]{
			\includegraphics[#7]{\defaultimagefolder#6}}
		\setcaptionmargincm{\captionlrmargin}
		\ifx\hfuzz#8\hfuzz
			\vspace{\captionlessmarginimage cm}
		\else
			\caption{#8 #1}
		\fi
	\end{figure}
	\setcaptionmargincm{\captionlrmargin}
	\vspace{\marginbottomimages cm}
}

\newcommand{\inserthexaimage}[9][]{
	% Insertar una imagen con 6 imágenes, igual propiedades
	% 	#1	Label (opcional)
	%	#2	Dirección de la imagen 1
	%	#3	Dirección de la imagen 2
	%	#4	Dirección de la imagen 3
	%	#5	Dirección de la imagen 4
	%	#6	Dirección de la imagen 5
	%	#7	Dirección de la imagen 6
	%	#8	Propiedades de las imágenes
	%	#9	Leyenda de la imagen global
	\emptyvarerr{\inserthexaimage}{#2}{Direccion de la imagen 1 no definida}
	\emptyvarerr{\inserthexaimage}{#3}{Direccion de la imagen 2 no definida}
	\emptyvarerr{\inserthexaimage}{#4}{Direccion de la imagen 3 no definida}
	\emptyvarerr{\inserthexaimage}{#5}{Direccion de la imagen 4 no definida}
	\emptyvarerr{\inserthexaimage}{#6}{Direccion de la imagen 5 no definida}
	\emptyvarerr{\inserthexaimage}{#7}{Direccion de la imagen 6 no definida}
	\emptyvarerr{\inserthexaimage}{#8}{Propiedades de las imagenes no definidas}
	\vspace{\margintopimages cm}
	\captionsetup{margin=0.45cm}
	\begin{figure}[H] \centering
		\subfloat[]{
			\includegraphics[#8]{\defaultimagefolder#2}}
		\hspace{0.1cm}
		\subfloat[]{
			\includegraphics[#8]{\defaultimagefolder#3}}
		\hspace{0.1cm}
		\subfloat[]{
			\includegraphics[#8]{\defaultimagefolder#4}}
		\hspace{0.1cm}
		\subfloat[]{
			\includegraphics[#8]{\defaultimagefolder#5}}
		\hspace{0.1cm}
		\subfloat[]{
			\includegraphics[#8]{\defaultimagefolder#6}}
		\hspace{0.1cm}
		\subfloat[]{
			\includegraphics[#8]{\defaultimagefolder#7}}
		\setcaptionmargincm{\captionlrmargin}
		\ifx\hfuzz#9\hfuzz
			\vspace{\captionlessmarginimage cm}
		\else
			\caption{#9 #1}
		\fi
	\end{figure}
	\setcaptionmargincm{\captionlrmargin}
	\vspace{\marginbottomimages cm}
}

\newcommand{\insertimageleft}[5][]{
	% Insertar una imagen a la izquierda
	% 	#1	Label (opcional)
	%	#2	Dirección de la imagen
	%	#3	Ancho de la imagen (en textwidth)
	%	#4	Altura en líneas de la imagen
	%	#5	Leyenda de la imagen
	\emptyvarerr{\insertimageleft}{#2}{Direccion de la imagen no definida}
	\emptyvarerr{\insertimageleft}{#3}{Ancho de la imagen no defindo}
	\emptyvarerr{\insertimageleft}{#4}{Altura en lineas de la imagen no definida}
	\begin{wrapfigure}{L}{#3\textwidth}
		\setcaptionmargincm{0}
		\ifthenelse{\equal{\figurecaptiontop}{true}}{}{
		\vspace{\marginfloatimages pt}}
		\centering
		\includegraphics[width=\linewidth]{#2}
		\ifx\hfuzz#5\hfuzz
			\vspace{\captionlessmarginimage cm}
		\else
			\caption{#5 #1}
		\fi
		\setcaptionmargincm{\captionlrmargin}
	\end{wrapfigure}
}

\newcommand{\insertimageright}[5][]{
	% Insertar una imagen a la derecha
	% 	#1	Label (opcional)
	%	#2	Dirección de la imagen
	%	#3	Ancho de la imagen (en textwidth)
	%	#4	Altura en líneas de la imagen
	%	#5	Leyenda de la imagen
	\emptyvarerr{\insertimageright}{#2}{Direccion de la imagen no definida}
	\emptyvarerr{\insertimageright}{#3}{Ancho de la imagen no defindo}
	\emptyvarerr{\insertimageright}{#4}{Altura en lineas de la imagen no definida}
	\begin{wrapfigure}{r}{#3\textwidth}
		\setcaptionmargincm{0}
		\ifthenelse{\equal{\figurecaptiontop}{true}}{}{
			\vspace{\marginfloatimages pt}}
		\centering
		\includegraphics[width=1.0\textwidth]{#2}
		\ifx\hfuzz#5\hfuzz
			\vspace{\captionlessmarginimage cm}
		\else
			\caption{#5 #1}
		\fi
		\setcaptionmargincm{\captionlrmargin}
	\end{wrapfigure}
}

% Inserta un texto entre comillas
\newcommand{\quotes}[1]{``#1''}

% Inserta un email con un link cliqueable
\newcommand{\insertemail}[1]{\href{mailto:#1}{\texttt{#1}}}

% IMPORTACIÓN DE ENTORNOS
% Template:     Informe/Reporte LaTeX
% Documento:    Definición de entornos
% Versión:      4.0.4 (20/06/2017)
% Codificación: UTF-8
%
% Autor: Pablo Pizarro R.
%        Facultad de Ciencias Físicas y Matemáticas
%        Universidad de Chile
%        pablo.pizarro@ing.uchile.cl, ppizarror.com
%
% Manual template: [http://ppizarror.com/Template-Informe/]
% Licencia MIT:    [https://opensource.org/licenses/MIT/]

% Crea una sección de referencias solo para bibtex
\newenvironment{references}{
	\ifthenelse{\equal{\stylecitereferences}{bibtex}}{
	}{
		\throwerror{}{Solo se puede usar entorno references con estilo citas \noexpand\stylecitereferences=bibtex}
	}
	\begingroup
	% Se configura las referencias como una sección
	\ifthenelse{\equal{\referencenumsection}{true}}{
		\section{\namereferences}
	}{
		\sectionanum{\namereferences}
	}
	\renewcommand{\section}[2]{}
	\begin{thebibliography}{99}
	}
	{ 
	\end{thebibliography}
	\endgroup
}

% Crea una sección de resumen
\newenvironment{resumen}{
	% Tipo de título para abstract
	\sectionfont{\color{\titlecolor} \fontsizetitle \styletitle \selectfont}
	
	% Inserta un título sin número, sin cabecera y sin aparecer en el índice, para que aparezca en el índice utilizar la función \sectionanumheadless
	\sectionanumnoiheadless{\nameabstract}}{
	
	% Salta de página si está imprimiendo por ambas caras
	\ifthenelse{\equal{\addemptypagetwosides}{true}}{
		\checkoddpage
		\ifoddpage
			\newpage
			\null
			\thispagestyle{empty}
			\newpage
			\addtocounter{page}{-1}
		\else
		\fi
	}{}
}

% Columna centrada en tablas
\newcolumntype{P}[1]{
	>{\centering\arraybackslash}p{#1}
}

% IMPORTACIÓN DE ESTILOS
% Template:     Informe/Reporte LaTeX
% Documento:    Definición de estilos
% Versión:      4.5.2 (04/09/2017)
% Codificación: UTF-8
%
% Autor: Pablo Pizarro R.
%        Facultad de Ciencias Físicas y Matemáticas
%        Universidad de Chile
%        pablo.pizarro@ing.uchile.cl, ppizarror.com
%
% Manual template: [http://latex.ppizarror.com/Template-Informe/]
% Licencia MIT:    [https://opensource.org/licenses/MIT/]

% Definición de colores
\definecolor{backcolour}{rgb}{0.95, 0.95, 0.92}
\definecolor{codegray}{rgb}{0.5, 0.5, 0.5}
\definecolor{codegreen}{rgb}{0, 0.6, 0}
\definecolor{codepurple}{rgb}{0.58, 0, 0.82}
\definecolor{dgray}{RGB}{104, 108, 113}
\definecolor{dkgreen}{rgb}{0, 0.6, 0}
\definecolor{gray}{rgb}{0.5, 0.5, 0.5}
\definecolor{lbrown}{RGB}{255, 252, 249}
\definecolor{lgray}{RGB}{180, 180, 180}
\definecolor{lyellow}{rgb}{1.0, 1.0, 0.88}
\definecolor{mauve}{rgb}{0.58, 0, 0.82}
\definecolor{mygray}{rgb}{0.5, 0.5, 0.5}
\definecolor{mygreen}{rgb}{0, 0.6, 0}
\definecolor{mylilas}{RGB}{170, 55, 241}
\definecolor{ocre}{RGB}{243, 102, 25}

% Estilo de códigos
\ifthenelse{\equal{\codecaptiontop}{true}}{ % Leyenda arriba

	% Estilo de lenguaje C
	\lstdefinestyle{C}{
		language=C,
		aboveskip=3mm,
		backgroundcolor=\color{backcolour},
		basicstyle={\small\ttfamily},
		belowskip=3mm,
		breakatwhitespace=false,
		breaklines=true,
		captionpos=t,
		columns=flexible,
		commentstyle=\color{mygreen},
		keepspaces=true,
		keywordstyle=\color{magenta},
		numbers=left,
		numbersep=5pt,
		numberstyle=\tiny\color{mygray},
		showspaces=false,
		showstringspaces=false,
		showtabs=false,
		stepnumber=1,
		stringstyle=\color{mauve},
		tabsize=3
	}
	
	% Estilo de lenguaje Java
	\lstdefinestyle{Java}{
		language=Java,
		aboveskip=3mm,
		backgroundcolor=\color{backcolour},
		basicstyle={\small\ttfamily},
		belowskip=3mm,
		breakatwhitespace=true,
		breaklines=true,
		captionpos=t,
		columns=flexible,
		commentstyle=\color{dkgreen},
		keepspaces=true,
		keywordstyle=\color{blue},
		numbers=left,
		numbersep=5pt,
		numberstyle=\tiny\color{gray},
		showstringspaces=false,
		stepnumber=1,
		stringstyle=\color{mauve},
		tabsize=3
	}
	
	% Estilo de lenguaje Matlab
	\lstdefinestyle{Matlab}{
		language=Matlab,
		aboveskip=3mm,
		backgroundcolor=\color{backcolour},
		basicstyle={\small\ttfamily},
		belowskip=3mm,
		breaklines=true,
		breaklines=true,
		captionpos=t,
		columns=flexible,
		commentstyle=\color{mygreen},
		emph=[1]{for,end,break},emphstyle=[1]\color{red},
		identifierstyle=\color{black},
		keepspaces=true,
		keywordstyle=\color{blue},
		morekeywords={matlab2tikz},
		numbers=left,
		numbersep=5pt,
		numberstyle=\tiny\color{gray},
		showstringspaces=false,
		stepnumber=1,
		stringstyle=\color{mylilas},
		tabsize=3
	}
	
	% Estilo de lenguaje Python
	\lstdefinestyle{Python}{
		language=Python,
		aboveskip=3mm,
		backgroundcolor=\color{backcolour},
		basicstyle={\small\ttfamily},
		belowskip=3mm,
		breakatwhitespace=false,
		breaklines=true,
		captionpos=t,
		columns=flexible,
		commentstyle=\color{codegreen},
		keepspaces=true,
		keywordstyle=\color{magenta},
		numbers=left,
		numbersep=5pt,
		numberstyle=\tiny\color{codegray},
		showspaces=false,
		showstringspaces=false,
		showtabs=false,
		stepnumber=1,
		stringstyle=\color{codepurple},
		tabsize=3
	}
}{ % Leyenda abajo

	% Estilo de lenguaje C
	\lstdefinestyle{C}{
		language=C,
		aboveskip=3mm,
		backgroundcolor=\color{backcolour},
		basicstyle={\small\ttfamily},
		belowskip=3mm,
		breakatwhitespace=false,
		breaklines=true,
		captionpos=b,
		columns=flexible,
		commentstyle=\color{mygreen},
		keepspaces=true,
		keywordstyle=\color{magenta},
		numbers=left,
		numbersep=5pt,
		numberstyle=\tiny\color{mygray},
		showspaces=false,
		showstringspaces=false,
		showtabs=false,
		stepnumber=1,
		stringstyle=\color{mauve},
		tabsize=3
	}
	
	% Estilo de lenguaje Java
	\lstdefinestyle{Java}{
		language=Java,
		aboveskip=3mm,
		backgroundcolor=\color{backcolour},
		basicstyle={\small\ttfamily},
		belowskip=3mm,
		breakatwhitespace=true,
		breaklines=true,
		captionpos=b,
		columns=flexible,
		commentstyle=\color{dkgreen},
		keepspaces=true,
		keywordstyle=\color{blue},
		numbers=left,
		numbersep=5pt,
		numberstyle=\tiny\color{gray},
		showstringspaces=false,
		stepnumber=1,
		stringstyle=\color{mauve},
		tabsize=3
	}
	
	% Estilo de lenguaje Matlab
	\lstdefinestyle{Matlab}{
		language=Matlab,
		aboveskip=3mm,
		backgroundcolor=\color{backcolour},
		basicstyle={\small\ttfamily},
		belowskip=3mm,
		breaklines=true,
		breaklines=true,
		captionpos=b,
		columns=flexible,
		commentstyle=\color{mygreen},
		emph=[1]{for,end,break},emphstyle=[1]\color{red},
		identifierstyle=\color{black},
		keepspaces=true,
		keywordstyle=\color{blue},
		morekeywords={matlab2tikz},
		numbers=left,
		numbersep=5pt,
		numberstyle=\tiny\color{gray},
		showstringspaces=false,
		stepnumber=1,
		stringstyle=\color{mylilas},
		tabsize=3
	}
	
	% Estilo de lenguaje Python
	\lstdefinestyle{Python}{
		language=Python,
		aboveskip=3mm,
		backgroundcolor=\color{backcolour},
		basicstyle={\small\ttfamily},
		belowskip=3mm,
		breakatwhitespace=false,
		breaklines=true,
		captionpos=b,
		columns=flexible,
		commentstyle=\color{codegreen},
		keepspaces=true,
		keywordstyle=\color{magenta},
		numbers=left,
		numbersep=5pt,
		numberstyle=\tiny\color{codegray},
		showspaces=false,
		showstringspaces=false,
		showtabs=false,
		stepnumber=1,
		stringstyle=\color{codepurple},
		tabsize=3
	}	
}

% Estilo de enumeración en griego
\RequirePackage{enumitem}
\makeatletter
\def\greek#1{\expandafter\@greek\csname c@#1\endcsname}
\def\Greek#1{\expandafter\@Greek\csname c@#1\endcsname}
\def\@greek#1{\ifcase#1
	\or $\alpha$%
	\or $\beta$%
	\or $\gamma$%
	\or $\delta$%
	\or $\epsilon$%
	\or $\zeta$%
	\or $\eta$%
	\or $\theta$%
	\or $\iota$%
	\or $\kappa$%
	\or $\lambda$%
	\or $\mu$%
	\or $\nu$%
	\or $\xi$%
	\or $o$%
	\or $\pi$%
	\or $\rho$%
	\or $\sigma$%
	\or $\tau$%
	\or $\upsilon$%
	\or $\phi$%
	\or $\chi$%
	\or $\psi$%
	\or $\omega$%
	\fi}
\def\@Greek#1{\ifcase#1
	\or $\mathrm{A}$%
	\or $\mathrm{B}$%
	\or $\Gamma$%
	\or $\Delta$%
	\or $\mathrm{E}$%
	\or $\mathrm{Z}$%
	\or $\mathrm{H}$%
	\or $\Theta$%
	\or $\mathrm{I}$%
	\or $\mathrm{K}$%
	\or $\Lambda$%
	\or $\mathrm{M}$%
	\or $\mathrm{N}$%
	\or $\Xi$%
	\or $\mathrm{O}$%
	\or $\Pi$%
	\or $\mathrm{P}$%
	\or $\Sigma$%
	\or $\mathrm{T}$%
	\or $\mathrm{Y}$%
	\or $\Phi$%
	\or $\mathrm{X}$%
	\or $\Psi$%
	\or $\Omega$%
	\fi}
\makeatother
\AddEnumerateCounter{\greek}{\@greek}{24}
\AddEnumerateCounter{\Greek}{\@Greek}{12}

% Columna centrada en tablas
\newcolumntype{P}[1]{
	>{\centering\arraybackslash}p{#1}
}


% CONFIGURACIÓN INICIAL DEL DOCUMENTO
% Template:     Informe/Reporte LaTeX
% Documento:    Configuración inicial del template
% Versión:      4.2.7 (09/07/2017)
% Codificación: UTF-8
%
% Autor: Pablo Pizarro R.
%        Facultad de Ciencias Físicas y Matemáticas
%        Universidad de Chile
%        pablo.pizarro@ing.uchile.cl, ppizarror.com
%
% Manual template: [http://ppizarror.com/Template-Informe/]
% Licencia MIT:    [https://opensource.org/licenses/MIT/]

% Se revisa si las variables no han sido borradas
\checkvardefined{\nombredelinforme}
\checkvardefined{\temaatratar}
\checkvardefined{\autordeldocumento}
\checkvardefined{\nombredelcurso}
\checkvardefined{\codigodelcurso}
\checkvardefined{\nombreuniversidad}
\checkvardefined{\nombrefacultad}
\checkvardefined{\departamentouniversidad}
\checkvardefined{\localizacionuniversidad}

% Se añade \xspace a las variables
\makeatletter
	\g@addto@macro\nombredelinforme\xspace
	\g@addto@macro\temaatratar\xspace
	\g@addto@macro\autordeldocumento\xspace
	\g@addto@macro\nombredelcurso\xspace
	\g@addto@macro\codigodelcurso\xspace
	\g@addto@macro\nombreuniversidad\xspace
	\g@addto@macro\nombrefacultad\xspace
	\g@addto@macro\departamentouniversidad\xspace
	\g@addto@macro\localizacionuniversidad\xspace
\makeatother

% Se crean variables si se borraron
\ifthenelse{\isundefined{\tablaintegrantes}}{
	\errmessage{LaTeX Warning: Se borro la variable \noexpand\tablaintegrantes, creando una vacia}
	\def\tablaintegrantes {}
}{}

% Se define metadata del pdf
\hypersetup{
	bookmarksopen={\cfgpdfbookmarkopen},
	bookmarksopenlevel={\cfgbookmarksopenlevel},
	bookmarkstype={toc},
	pdfauthor={\autordeldocumento},
	pdfcenterwindow={\cfgpdfcenterwindow},
	pdfcreator={LaTeX, pdfLaTeX},
	pdfdisplaydoctitle={\cfgpdfdisplaydoctitle},
	pdffitwindow={\cfgpdffitwindow},
	pdfinfo={
		Curso.Codigo={\codigodelcurso},
		Curso.Nombre={\nombredelcurso},
		Documento.Autor={\autordeldocumento},
		Documento.Tema={\temaatratar},
		Documento.Titulo={\nombredelinforme},
		Template.Autor.Alias={ppizarror},
		Template.Autor.Email={pablo.pizarro@ing.uchile.cl},
		Template.Autor.Nombre={Pablo Pizarro R.},
		Template.Autor.Web={http://ppizarror.com/},
		Template.Codificacion={UTF-8},
		Template.Fecha={09/07/2017},
		Template.Licencia.Tipo={MIT},
		Template.Licencia.Web={https://opensource.org/licenses/MIT/},
		Template.Nombre={Template-Informe},
		Template.Tipo={Normal},
		Template.Version.Dev={4.2.7-1-N},
		Template.Version.Hash={6CC0BE369E23588566C56F922A7FF0C1},
		Template.Version.Release={4.2.7},
		Template.Web.Dev={https://github.com/ppizarror/Template-Informe/},
		Template.Web.Manual={http://ppizarror.com/Template-Informe/},
		Universidad.Departamento={\departamentouniversidad},
		Universidad.Nombre={\nombreuniversidad},
		Universidad.Ubicacion={\localizacionuniversidad},
	},
	pdfkeywords={\nombreuniversidad, \codigodelcurso \nombredelcurso, \localizacionuniversidad},
	pdfpagemode={UseOutlines},
	pdfproducer={Template-Informe v4.2.7 | (Pablo Pizarro R.) ppizarror.com},
	pdfstartpage={1},
	pdfstartview={FitH},
	pdfsubject={\temaatratar},
	pdftitle={\nombredelinforme},
	pdftoolbar={\cfgpdftoolbar},
	pdfview={FitH}
}

% Establece la carpeta de imágenes por defecto
\graphicspath{{./\defaultimagefolder}}

% Definición de valores e dimensiones
\setlength{\headheight}{64pt} % Tamaño de la cabecera sin fancyhdr
\setcounter{MaxMatrixCols}{20} % Número máximo de columnas en matrices
\newcounter{templatepagecounter} % Administra números de páginas
\setlength{\footnotemargin}{3mm} % Margen del footnote
\renewcommand{\baselinestretch}{\defaultinterline} % Ajuste del entrelineado

% Configuración de los colores
\color{\maintextcolor} % Color principal
\arrayrulecolor{\tablelinecolor} % Color de las lineas de las tablas
\sethlcolor{\highlightcolor} % Color del subrayado por defecto
\ifthenelse{\equal{\showborderonlinks}{true}}{
	\hypersetup{
		% Color de links con borde
		citebordercolor=\citecolor,
		linkbordercolor=\linkcolor,
		urlbordercolor=\urlcolor
	}
}{
	\hypersetup{
		% Color de links sin borde
		hidelinks,
		colorlinks=true,
		citecolor=\citecolor,
		linkcolor=\linkcolor,
		urlcolor=\urlcolor
	}
}

% Configuración de las leyendas
\setcaptionmargincm{\captionlrmargin} % Márgenes de las leyendas por defecto
\ifthenelse{\equal{\captiontextbold}{true}}{ % Texto en negrita en etiquetas
	\renewcommand{\captiontextbold}{bf}}{
	\renewcommand{\captiontextbold}{}
}
\captionsetup{ % Se configura el texto de los caption
	labelfont={color=\captioncolor, \captiontextbold},
	textfont={color=\captiontextcolor},
	singlelinecheck=on
}
\floatsetup[figure]{ % Configuración de márgenes en las figuras
	captionskip=\captiontbmarginfigure pt
}
\floatsetup[table]{ % Configuración de márgenes en las tablas
	captionskip=\captiontbmargintable pt
}
\ifthenelse{\equal{\figurecaptiontop}{true}}{ % Caption superior en figuras
	\floatsetup[figure]{position=above}}{
}
\ifthenelse{\equal{\tablecaptiontop}{true}}{ % Caption superior en tablas
	\floatsetup[table]{position=top}
	}{
	\floatsetup[table]{position=bottom}
}
\ifthenelse{\equal{\captionalignment}{justified}}{ % Leyenda justificada
	\captionsetup{
		justification=justified,
		format=plain
	}
}{
	\ifthenelse{\equal{\captionalignment}{centered}}{ % Leyenda centrada
		\captionsetup{
			justification=centering
		}
	}{
		\ifthenelse{\equal{\captionalignment}{left}}{ % Leyenda izq
			\captionsetup{
				justification=raggedright,
				singlelinecheck=false
			}
		}{
			\ifthenelse{\equal{\captionalignment}{right}}{ % Leyenda der
				\captionsetup{
					justification=raggedleft,
					singlelinecheck=false
				}
			}{
				\throwbadconfig{Posicion de leyendas desconocida}{\captionalignment}{justified,centered,left,right}
			}
		}
	}
}

% Configuración de referencias y citas
\ifthenelse{\equal{\stylecitereferences}{natbib}}{
	\bibliographystyle{apa}
	\setlength{\bibsep}{\natbibrefsep pt}
}{
	\ifthenelse{\equal{\stylecitereferences}{apacite}}{
		\bibliographystyle{apacite}
		\setlength{\bibitemsep}{\apaciterefsep pt}
	}{
		\ifthenelse{\equal{\stylecitereferences}{bibtex}}{
			\bibliographystyle{apa}
			\newlength{\bibitemsep}\setlength{\bibitemsep}{.2\baselineskip plus .05\baselineskip minus .05\baselineskip}
			\newlength{\bibparskip}\setlength{\bibparskip}{0pt}
			\let\oldthebibliography\thebibliography
			\renewcommand\thebibliography[1]{
				\oldthebibliography{#1}
				\setlength{\parskip}{\bibitemsep}
				\setlength{\itemsep}{\bibparskip}
			}
			\setlength{\bibitemsep}{\bibtexrefsep pt}
		}{
			\throwbadconfig{Estilo citas desconocido}{\stylecitereferences}{bibtex,apacite,natbib}
		}
	}
}
\makeatletter
\ifthenelse{\equal{\twocolumnreferences}{true}}{
	% Referencias en 2 columnas
	\renewenvironment{thebibliography}[1]
	{\begin{multicols}{2}[\section*{\refname}]
		\@mkboth{\MakeUppercase\refname}{\MakeUppercase\refname}
		\list{\@biblabel{\@arabic\c@enumiv}}
		{\settowidth\labelwidth{\@biblabel{#1}}
			\leftmargin\labelwidth
			\advance\leftmargin\labelsep
			\@openbib@code
			\usecounter{enumiv}
			\let\p@enumiv\@empty
			\renewcommand\theenumiv{\@arabic\c@enumiv}}
		\sloppy
		\clubpenalty 4000
		\@clubpenalty \clubpenalty
		\widowpenalty 4000
		\sfcode`\.\@m}
		{\def\@noitemerr
		{\@latex@warning{Ambiente `thebibliography' no definido}}
		\endlist\end{multicols}}}{}
\makeatother

% Configuración anexo
\patchcmd{\appendices}{\quad}{\sectionappendixlastchar\quad}{}{}

% Configuración de códigos fuente
\lstset{
	extendedchars=true,
	keepspaces=true,
	columns=flexible,
	literate=
	{á}{{\'a}}1 {é}{{\'e}}1 {í}{{\'i}}1 {ó}{{\'o}}1 {ú}{{\'u}}1
	{Á}{{\'A}}1 {É}{{\'E}}1 {Í}{{\'I}}1 {Ó}{{\'O}}1 {Ú}{{\'U}}1
	{à}{{\`a}}1 {è}{{\`e}}1 {ì}{{\`i}}1 {ò}{{\`o}}1 {ù}{{\`u}}1
	{À}{{\`A}}1 {È}{{\'E}}1 {Ì}{{\`I}}1 {Ò}{{\`O}}1 {Ù}{{\`U}}1
	{ä}{{\"a}}1 {ë}{{\"e}}1 {ï}{{\"i}}1 {ö}{{\"o}}1 {ü}{{\"u}}1
	{Ä}{{\"A}}1 {Ë}{{\"E}}1 {Ï}{{\"I}}1 {Ö}{{\"O}}1 {Ü}{{\"U}}1
	{â}{{\^a}}1 {ê}{{\^e}}1 {î}{{\^i}}1 {ô}{{\^o}}1 {û}{{\^u}}1
	{Â}{{\^A}}1 {Ê}{{\^E}}1 {Î}{{\^I}}1 {Ô}{{\^O}}1 {Û}{{\^U}}1
	{œ}{{\oe}}1 {Œ}{{\OE}}1 {æ}{{\ae}}1 {Æ}{{\AE}}1 {ß}{{\ss}}1
	{ű}{{\H{u}}}1 {Ű}{{\H{U}}}1 {ő}{{\H{o}}}1 {Ő}{{\H{O}}}1
	{ç}{{\c c}}1 {Ç}{{\c C}}1 {ø}{{\o}}1 {å}{{\r a}}1 {Å}{{\r A}}1
	{€}{{\EUR}}1 {£}{{\pounds}}1
}

% Configuración de hbox y vbox
\hfuzz=200pt \vfuzz=200pt
\hbadness=\maxdimen \vbadness=\maxdimen

% Se activa el word-wrap para textos con \texttt{}
\ttfamily \hyphenchar\the\font=`\-
	
% Se define el tipo de texto de los url
\urlstyle{tt}

% Se activa el modo estricto de revisión de números de página
\strictpagecheck


% INICIO DE LAS PÁGINAS
\begin{document}

% PORTADA
% Template:     Informe/Reporte LaTeX
% Documento:    Definición de portadas
% Versión:      4.5.5 (04/09/2017)
% Codificación: UTF-8
%
% Autor: Pablo Pizarro R.
%        Facultad de Ciencias Físicas y Matemáticas
%        Universidad de Chile
%        pablo.pizarro@ing.uchile.cl, ppizarror.com
%
% Manual template: [http://latex.ppizarror.com/Template-Informe/]
% Licencia MIT:    [https://opensource.org/licenses/MIT/]

% Se escribe el header de la portada
\newpage
\renewcommand{\thepage}{\nameportraitpage}

% Estilos de portada
\ifthenelse{\equal{\portraitstyle}{style1}}{
	\setpagemargincm{\pagemarginleft}{\firstpagemargintop}{\pagemarginright}{\pagemarginbottom}
	\pagestyle{fancy} \fancyhf{}
	\fancyhead[L]{
		\nombreuniversidad \\ \nombrefacultad \\ \departamentouniversidad \\ \vspace{-0.43cm}
	}
	\fancyhead[R]{
		\includegraphics[scale=\imagendepartamentoescala]{\imagendepartamento}
	}

	~ \\
	\vfill
	\begin{center}
		{\noindent \LARGE{\nombredelcurso} \vspace{0.3cm}} ~ \\
		\vspace*{1.5cm}
		\textcolor{\portraittitlecolor}{
			{\centering \noindent \Huge{\titulodelinforme} \vspace{0.3cm}} ~ \\
			{\noindent \large{\temaatratar}}
		}
	\end{center}

	\vfill
	\noindent
	\begin{minipage}{1.0\textwidth}
		\begin{flushright}
			\tablaintegrantes
		\end{flushright}
	\end{minipage}
}{
\ifthenelse{\equal{\portraitstyle}{style2}}{
	\setpagemargincm{\pagemarginleft}{\firstpagemargintop}{\pagemarginright}{\pagemarginbottom}
	\pagestyle{fancy} \fancyhf{}
	\fancyhead[L]{
		\nombreuniversidad \\ \nombrefacultad \\ \departamentouniversidad \\ \vspace{-0.43cm}
	}
	\fancyhead[R]{
		\includegraphics[scale=\imagendepartamentoescala]{\imagendepartamento}
	}
	
	~ \\
	\vfill
	\begin{center}
		\textcolor{\portraittitlecolor}{
			{\noindent \Huge{\titulodelinforme} \vspace{0.5cm}} ~ \\
			{\noindent \large{\temaatratar}}
		}
	\end{center}
	
	\vfill
	\noindent
	\begin{minipage}{1.0\textwidth}
		\begin{flushright}
			\tablaintegrantes
		\end{flushright}
	\end{minipage}
}{
\ifthenelse{\equal{\portraitstyle}{style3}}{
	\setpagemargincm{\pagemarginleft}{\firstpagemargintop}{\pagemarginright}{\pagemarginbottom}
	\pagestyle{fancy} \fancyhf{}
	\fancyhead[L]{
		\nombreuniversidad \\ \nombrefacultad \\ \departamentouniversidad \\ \vspace{-0.43cm}
	}
	\fancyhead[R]{
		\includegraphics[scale=\imagendepartamentoescala]{\imagendepartamento}
	}

	~ \\
	\vfill
	\begin{center}
		\vspace*{-1.0cm}
		{\noindent \huge{\nombredelcurso} \vspace{0.3cm}} ~ \\
		{\noindent \large{Código del curso: \codigodelcurso}} ~ \\
		\vspace*{1.8cm}
		\textcolor{\portraittitlecolor}{
			{\noindent \Huge{\titulodelinforme} \vspace{0.3cm}} ~ \\
			{\noindent \large{\temaatratar}}
		}
	\end{center}

	\vfill
	\noindent
	\begin{minipage}{1.0\textwidth}
		\begin{flushright}
			\tablaintegrantes
		\end{flushright}
	\end{minipage}
}{
\ifthenelse{\equal{\portraitstyle}{style4}}{
	\setpagemargincm{\pagemarginleft}{\pagemargintop}{\pagemarginright}{\pagemarginbottom}
	\thispagestyle{empty}
	
	\begin{wrapfigure}{l}{0.3\textwidth}
		\vspace{-0.69cm}
		\noindent \hspace{-1.10cm} \includegraphics[scale=1.35]{departamentos/fcfm2}
	\end{wrapfigure}
	\hspace*{0.3cm}
	\noindent \textsc{\color{red} \hspace{-2.2cm} \departamentouniversidad} \\
	\hspace*{0.3cm}
	\noindent \textsc{\color{dgray} \hspace{-1.6cm} \nombrefacultad} \\
	\hspace*{0.3cm}
	\noindent \textsc{\color{dgray} \hspace{-1.6cm} \nombreuniversidad} \\
	\hspace*{0.3cm}
	\noindent \textsc{\color{dgray} \hspace{-1.6cm} \codigodelcurso \nombredelcurso}\\
	
	\vfill
	\begin{center}
		\vspace*{0.5cm}
		{\color{dgray} \Large \textbf{\MakeUppercase{\temaatratar}}} \\
		\noindent \rule{\linewidth}{0.4mm} ~ \\
		\Huge \textup \bfseries \textsc{\textcolor{\portraittitlecolor}{\titulodelinforme}} \\
		\noindent \rule{\linewidth}{0.4mm} ~ \\
	\end{center}
	\begin{minipage}{.5\textwidth}
		~
	\end{minipage}
	
	\vfill
	\begin{minipage}{1.0\textwidth}
		\begin{flushright}
			\noindent \tablaintegrantes
		\end{flushright}
	\end{minipage}
}{
\ifthenelse{\equal{\portraitstyle}{style5}}{
	\setpagemargincm{\pagemarginleft}{\pagemargintop}{\pagemarginright}{\pagemarginbottom}
	\thispagestyle{empty}
	
	\vspace*{-1.5cm}
	\noindent \includegraphics[width=1.75cm]{departamentos/uchile2}
	\hspace*{-0.15cm}
	\begin{tabular}{l}
		\small \scshape{\MakeUppercase{\nombreuniversidad}} \\
		\small \scshape{\MakeUppercase{\nombrefacultad}} \\
		\small \scshape{\MakeUppercase{\departamentouniversidad}} \\
		\vspace*{1.25cm}\mbox{}
	\end{tabular}

	\vfill
	\begin{center}
		{\fontsize{22pt}{10pt} \selectfont 
		\noindent \textcolor{\portraittitlecolor}{\titulodelinforme} \vspace*{0.35cm}} \\
		{\noindent \fontsize{10pt}{5pt} \selectfont  \textcolor{\portraittitlecolor}{\codigodelcurso\ - \nombredelcurso}} \\
		\vspace*{3cm}
	\end{center}

	\vfill
	\noindent
	\begin{minipage}{1.0\textwidth}
		\begin{flushright}
			\tablaintegrantes
		\end{flushright}
	\end{minipage}
}{
\ifthenelse{\equal{\portraitstyle}{style6}}{
	\setpagemargincm{\pagemarginleft}{\pagemargintop}{\pagemarginright}{\pagemarginbottom}
	\thispagestyle{empty}
	
	\begin{center}
		\vspace*{-1.5cm}
		\includegraphics[scale=\imagendepartamentoescala]{\imagendepartamento}
		\hspace*{-0.15cm}
		\begin{tabular}{l}
			\vspace*{0.26cm}\mbox{} \\
			\small \scshape{\MakeUppercase{\nombreuniversidad}} \\
			\small \scshape{\MakeUppercase{\nombrefacultad}} \\
			\small \scshape{\MakeUppercase{\departamentouniversidad}} \\
			\vspace*{1.25cm}\mbox{}
		\end{tabular}
		
		\vfill
		\noindent \rule{\textwidth}{0.4mm} \\ \vspace{0.3cm}
		{\huge \textcolor{\portraittitlecolor}{\titulodelinforme} \vspace{0.2cm} \\}
		\noindent \rule{\textwidth}{0.4mm} \\ \vspace{0.40cm}
		{\large \textcolor{\portraittitlecolor}{\temaatratar} \\}
	\end{center}
	
	\vfill
	\noindent
	\begin{minipage}{1.0\textwidth}
		\begin{flushright}
			\scshape{\tablaintegrantes}
		\end{flushright}
	\end{minipage}
	
}{
\ifthenelse{\equal{\portraitstyle}{style7}}{
	\setpagemargincm{\pagemarginleft}{\pagemargintop}{\pagemarginright}{\pagemarginbottom}
	\thispagestyle{empty}
	
	\begin{center}
		\vspace*{-1.0cm}
		\begin{tabular}{c}
			\includegraphics[scale=\imagendepartamentoescala]{\imagendepartamento} \vspace{0.5cm} \\
			\small \scshape{\MakeUppercase{\nombreuniversidad}} \\
			\small \scshape{\MakeUppercase{\nombrefacultad}} \\
			\small \scshape{\MakeUppercase{\departamentouniversidad}}
		\end{tabular}
	\end{center}
	
	\vfill
	\begin{center}
		\noindent \rule{\textwidth}{0.4mm} \\ \vspace{0.3cm}
		{\huge \textcolor{\portraittitlecolor}{\titulodelinforme} \vspace{0.2cm} \\}
		\noindent \rule{\textwidth}{0.4mm} \\ \vspace{0.40cm}
		{\large \textcolor{\portraittitlecolor}{\temaatratar} \\}
	\end{center}
	
	\vfill
	\noindent
	\begin{minipage}{1.0\textwidth}
		\begin{flushright}
			\scshape{\tablaintegrantes}
		\end{flushright}
	\end{minipage}
}{
\ifthenelse{\equal{\portraitstyle}{style8}}{
	\setpagemargincm{\pagemarginleft}{\pagemargintop}{\pagemarginright}{\pagemarginbottom}
	\thispagestyle{empty}
	
	\includegraphics[width=1.5cm]{departamentos/uchile3} 
	\hspace{0cm}
	\begin{tabular}{l}
		\small \scshape{\MakeUppercase{\nombreuniversidad}} \\
		\small \scshape{\MakeUppercase{\nombrefacultad}} \\
		\small \scshape{\MakeUppercase{\departamentouniversidad}} \\
		\vspace*{1cm}\mbox{}
	\end{tabular}
	
	\vfill
	\begin{center}
		\fontsize{8mm}{9mm}\selectfont 
		\noindent \titulodelinforme \\
		\noindent \temaatratar
		
		\vspace*{1cm}
		\footnotesize{\codigodelcurso - \nombredelcurso} \\
		\vspace*{1.4cm}
	\end{center}
	
	\vfill
	\begin{center}
		\noindent \normalsize{\tablaintegrantes}
	\end{center}
}{
\ifthenelse{\equal{\portraitstyle}{style9}}{
	\setpagemargincm{\pagemarginleft}{\pagemargintop}{\pagemarginright}{\pagemarginbottom}
	\thispagestyle{empty}
	
	\noindent \includegraphics[scale=\imagendepartamentoescala]{\imagendepartamento}
	\vfill
	\begin{center}
		\noindent \rule{\textwidth}{0.4mm} \\ \vspace{0.3cm}
		{\huge \textcolor{\portraittitlecolor}{\titulodelinforme} \vspace{0.2cm} \\}
		\noindent \rule{\textwidth}{0.4mm} \\ \vspace{0.35cm}
		{\large \textcolor{\portraittitlecolor}{\temaatratar} \\}
	\end{center}
	
	\vfill
	\begin{center}
		\begin{tabular}{c}
			\small \scshape{\MakeUppercase{\nombreuniversidad}} \\
			\small \scshape{\MakeUppercase{\nombrefacultad}} \\
			\small \scshape{\MakeUppercase{\departamentouniversidad}}
		\end{tabular}
	\end{center}
	
	\vfill
	\begin{center}
		\indent \scshape{\tablaintegrantes}
	\end{center}
}{
\ifthenelse{\equal{\portraitstyle}{style10}}{
	\setpagemargincm{\pagemarginleft}{\pagemargintop}{\pagemarginright}{\pagemarginbottom}
	\thispagestyle{empty}
	
	\begin{center}
		\vspace*{-1.0cm}
		\scshape{\nombreuniversidad} \\
		\scshape{\nombrefacultad} \\
		\scshape{\departamentouniversidad}
	\end{center}
	
	\vfill
	\begin{center}
		{\setstretch{1.2} \fontsize{21pt}{22pt} \selectfont \textcolor{\portraittitlecolor}{\scshape{\titulodelinforme}} \vspace{0.5cm}} \\
		{\fontsize{13pt}{10pt} \selectfont \textcolor{\portraittitlecolor}{\scshape{\temaatratar}}}
	\end{center}
	
	\vfill
	\begin{center}
		\indent \tablaintegrantes
	\end{center}
}{
\ifthenelse{\equal{\portraitstyle}{style11}}{
	\setpagemargincm{\pagemarginleft}{\pagemargintop}{\pagemarginright}{\pagemarginbottom}
	\thispagestyle{empty}
	
	\begin{center}
		\vspace*{-1.0cm}
		\includegraphics[scale=\imagendepartamentoescala]{\imagendepartamento}
	\end{center}

	\vfill
	\begin{center}
		{\bf \Huge \scshape{\textcolor{\portraittitlecolor}{\titulodelinforme}} \vspace{0.3cm}} \\
		{\bf \Large \textcolor{\portraittitlecolor}{\temaatratar}}
	\end{center}
	
	\vfill
	\begin{flushright}
		\noindent \tablaintegrantes
	\end{flushright}

	\vspace{0.5cm}
	\noindent \rule{\textwidth}{0.4mm}
	\begin{center}
		\scshape{\nombreuniversidad, \nombrefacultad} \\
		\scshape{\departamentouniversidad}
	\end{center}
}{
\ifthenelse{\equal{\portraitstyle}{style12}}{
	\setpagemargincm{\pagemarginleft}{\pagemargintop}{\pagemarginright}{\pagemarginbottom}
	\thispagestyle{empty}
	
	\noindent
	\vspace*{-1.5cm}
	\begin{flushleft}
		\begin{minipage}{0.65\textwidth}
			{\fontsize{3.5mm}{0.5mm} \selectfont \noindent  \textsf{\nombreuniversidad, \nombrefacultad}} \\
			\noindent {\fontsize{3.0mm}{0.5mm} \selectfont \textsf{\departamentouniversidad} \vspace{-0.2cm}} \\
			\noindent \textcolor{lgray}{\rule{\textwidth}{0.3mm}}
		\end{minipage}
	\end{flushleft}
	\vspace*{-2.15cm}
	\begin{flushright}
		\begin{minipage}{0.3\textwidth}
			\noindent \includegraphics[width=1.0\textwidth]{\imagendepartamento}
		\end{minipage}
	\end{flushright}
	
	\vfill
	\begin{center}
		\begin{minipage}{0.9\textwidth}
			\begin{framed}
				\LARGE
				\vspace{1cm}
				\centering \textcolor{\portraittitlecolor}{\textbf{\titulodelinforme}}
				\vspace{1cm}
			\end{framed}
		\end{minipage}
	\end{center}
	
	\vfill
	\begin{flushright}
		\noindent \textsf{\tablaintegrantes}
	\end{flushright}
}{
\ifthenelse{\equal{\portraitstyle}{style13}}{
	\setpagemargincm{\pagemarginleft}{\pagemargintop}{\pagemarginright}{\pagemarginbottom}
	\thispagestyle{empty}
	
	\noindent
	\begin{flushleft}
		\vspace*{-1.0cm}
		\noindent \includegraphics[scale=\imagendepartamentoescala]{\imagendepartamento} \\
	\end{flushleft}
	
	\vfill
	{\bf \huge \noindent \textcolor{\portraittitlecolor}{\textsf{\MakeUppercase{\titulodelinforme}} \vspace*{0.05cm}}} \\
	{\bf \large \noindent \textcolor{\portraittitlecolor}{\textsf{\MakeUppercase{\temaatratar}}}} \\
	
	\vfill
	\begin{flushright}
		\noindent \textsf{\tablaintegrantes}
	\end{flushright}
}{
\ifthenelse{\equal{\portraitstyle}{style14}}{
	\setpagemargincm{\pagemarginleft}{\pagemargintop}{\pagemarginright}{\pagemarginbottom}
	\thispagestyle{empty}
	
	% Se chequean variables propias de la portada
	\checkextravarexist{\imagenextraportada}{}
	\checkextravarexist{\imagenextraportadaescala}{}
	
	\vspace*{-1.5cm}
	\noindent \begin{minipage}{0.8\textwidth}
		\noindent \begin{minipage}{0.22\textwidth}
			\includegraphics[scale=1.0]{departamentos/fcfm2} \\
		\end{minipage}
		\begin{minipage}{0.6\textwidth}
			\begin{flushleft}
				\textsc{
				\begin{tabular}{l}
					{\small \nombreuniversidad} \\
					{\small \nombrefacultad} \\
					{\small \departamentouniversidad}
				\end{tabular}
				}
			\end{flushleft}
		\end{minipage}
	\end{minipage}
	\noindent \begin{minipage}{0.2\textwidth}
		\begin{flushright}
			\ifthenelse{\isundefined{\imagenextraportada}}{}{
				\ifthenelse{\isundefined{\imagenextraportadaescala}}{}{
					\noindent \includegraphics[scale=\imagenextraportadaescala]{\imagenextraportada} \\
				}
			}
		\end{flushright}
	\end{minipage}
	
	\vfill
	\begin{center}
		{\fontsize{25pt}{15pt} \selectfont \textcolor{\portraittitlecolor}{\textbf{\titulodelinforme}} \vspace{0.7cm}} \\
		{\Large \textcolor{\portraittitlecolor}{\temaatratar}}
	\end{center}
	
	\vfill
	\begin{center}
		\noindent \tablaintegrantes
	\end{center}
}{
\ifthenelse{\equal{\portraitstyle}{style15}}{
	\setpagemargincm{\pagemarginleft}{\pagemargintop}{\pagemarginright}{\pagemarginbottom}
	
	\checkextravarexist{\imagenfondoportada}{Defina el fondo de la portada en el bloque INFORMACION DEL DOCUMENTO. Por defecto se uso la portada ubicada en ejemplos}
	\ifthenelse{\isundefined{\imagenfondoportada}}{
		\def\imagenfondoportada {ejemplos/background-1}
	}{}
	\checkextravarexist{\colorbloqueportada}{Defina el color del bloque del titulo de la portada en el bloque INFORMACION DEL DOCUMENTO. Por defecto se uso el color ocre}
	\ifthenelse{\isundefined{\imagenfondoportada}}{
		\def\imagenfondoportada {ejemplos/background-1}
	}{}
	\ifthenelse{\isundefined{\colorbloqueportada}}{
		\def\colorbloqueportada {ocre}
	}{}
	
	\begingroup
	\thispagestyle{empty}
	\begin{tikzpicture}[remember picture,overlay]
		\node[inner sep=0pt] (background) at (current page.center) {\includegraphics[width=\paperwidth]{\imagenfondoportada}};
		\draw (current page.center) node [fill=\colorbloqueportada!30!white,fill opacity=0.6,text opacity=1,inner sep=1cm]{\Huge\centering\bfseries\sffamily\parbox[c][][t]{\paperwidth}{
				\centering \textcolor{\portraittitlecolor}{\titulodelinforme} \\ [10pt]
				{\Large \temaatratar} \\ [25pt]
				{\huge \autordeldocumento}}};
	\end{tikzpicture}
	\vfill
	\endgroup
}{
\ifthenelse{\equal{\portraitstyle}{\bgtemplatetestcode}}{
	\setpagemargincm{\pagemarginleft}{\pagemargintop}{\pagemarginright}{\pagemarginbottom}
	\pagestyle{empty}
	\pagecolor{lbrown}
	
	\begin{center}
		\vspace*{-1.0cm}
		\scshape{\nombreuniversidad} \\
		\scshape{\nombrefacultad} \\
		\scshape{\departamentouniversidad}
	\end{center}
	
	~ \\
	\begin{center}
		\bgtemplatetestimg
	\end{center}

	\begin{center}
		\vspace*{-6cm}
		{\setstretch{1.2} \fontsize{25pt}{22pt} \selectfont \textcolor{\portraittitlecolor}{\scshape{\titulodelinforme}} \vspace{0.5cm}} \\
		{\fontsize{15pt}{10pt} \selectfont \textcolor{\portraittitlecolor}{\scshape{\temaatratar}}}
	\end{center}

	\vfill
	\begin{flushright}
		\noindent \tablaintegrantes
	\end{flushright}
	\newpage
	\pagecolor{white}
}{
	\throwbadconfigondoc{Estilo de portada incorrecto}{\portraitstyle}{style1 .. style15}}}}}}}}}}}}}}}}
}

% Añade una página en blanco al imprimir por las dos caras
\ifthenelse{\equal{\addemptypagetwosides}{true}}{
	\newpage
	\null
	\thispagestyle{empty}
	\renewcommand{\thepage}{}
	\newpage}{
}


% CONFIGURACIÓN DE PÁGINA Y ENCABEZADOS
% Template:     Informe/Reporte LaTeX
% Documento:    Configuración de página
% Versión:      3.0.0 (01/04/2017)
% Codificación: UTF-8
%
% Autor: Pablo Pizarro R.
%        Facultad de Ciencias Físicas y Matemáticas.
%        Universidad de Chile.
%        pablo.pizarro@ing.uchile.cl, ppizarror.com
%
% Sitio web del proyecto: [http://ppizarror.com/Template-Informe/]
% Licencia: MIT           [https://opensource.org/licenses/MIT]

\newpage
\pagenumbering{Roman}
\setcounter{page}{1}
\setcounter{footnote}{1}
\setpagemargincm{\defaultpagemarginleft}{\defaultpagemargintop}
{\defaultpagemarginright}{\defaultpagemarginbottom}
\def\arraystretch{\tablepadding} % Se ajusta el padding de las tablas
\renewcommand{\sectionmark}[1]{\markboth{#1}{}} % Se modifica el estilo del header
\renewcommand{\listfigurename}{\nomltfiguras} % Nombre del índice de figuras
\renewcommand{\listtablename}{\nomlttablas} % Nombre del índice de tablas
\renewcommand{\contentsname}{\nomltcontend} % Nombre del índice
\renewcommand{\lstlistlistingname}{\nomltcodfuente} % Nombre índice código fuente
\renewcommand{\tablename}{\nomltwtablas} % Nombre de la leyenda de las tablas
\renewcommand{\figurename}{\nomltwfigura} % Nombre de la leyenda de las figuras
\renewcommand{\lstlistingname}{\nomltsrc} % Nombre leyenda del código fuente
\pagestyle{fancy} \fancyhf{} % Se crean los headers y footers
\ifthenelse{\equal{\showheadertitle}{true}}{ % Header izq, nombre sección
	\fancyhead[L]{\nouppercase{\rightmark}}
}{}
\fancyhead[R]{\small \rm \nouppercase{\thepage}} % Header der, número página
\ifthenelse{\equal{\showfooter}{true}}{
	\fancyfoot[L]{
		\small \rm \textit{\nombredelinforme} % Footer izq, título del informe
	}
	\fancyfoot[R]{
		\small \rm \textit{\codigodelcurso \ \nombredelcurso} % Footer der, curso
	}
	\renewcommand{\footrulewidth}{0.5pt} % Ancho de la barra del footer
}{}
\renewcommand{\headrulewidth}{0.5pt} % Ancho de la barra del header
\sectionfont{
	\tipofuentetitulo \etipofuentetitulo \selectfont % Tipo de título para abstract
}

% RESUMEN O ABSTRACT
\begin{resumen}
	\lipsum[1] % Se incluye un párrafo de resumen, se puede borrar
\end{resumen}

% TABLA DE CONTENIDOS - ÍNDICE
% Template:     Informe/Reporte LaTeX
% Documento:    Índice
% Versión:      4.2.0 (07/07/2017)
% Codificación: UTF-8
%
% Autor: Pablo Pizarro R.
%        Facultad de Ciencias Físicas y Matemáticas
%        Universidad de Chile
%        pablo.pizarro@ing.uchile.cl, ppizarror.com
%
% Manual template: [http://ppizarror.com/Template-Informe/]
% Licencia MIT:    [https://opensource.org/licenses/MIT/]

\ifthenelse{\equal{\showindex}{true}}{
	% Crea nueva página y establece estilo de títulos
	\newpage
	\sectionfont{\color{\indextitlecolor} \fontsizetitlei \styletitlei \selectfont}
	
	% Índice de Contenidos
	\ifthenelse{\equal{\showindexofcontents}{true}}{\tableofcontents}{}
	
	% Lista de Figuras
	\iftotalfigures
		\ifthenelse{\equal{\showindexoffigures}{true}}{\listoffigures}{}
	\fi
	
	% Lista de Tablas
	\iftotaltables
		\ifthenelse{\equal{\showindexoftables}{true}}{\listoftables}{}
	\fi
	
	% Lista del Código Fuente
	\iftotallstlistings
		\ifthenelse{\equal{\showindexofcode}{true}}{\lstlistoflistings}{}
	\fi
	
	% Se añade una página en blanco
	\ifthenelse{\equal{\addemptypagetwosides}{true}}{
		\vfill
		\checkoddpage
		\ifoddpage
		\else
			\newpage
			\null
			\thispagestyle{empty}
			\newpage
			\addtocounter{page}{-1}
		\fi
	}{}
}{}
 % Índice, se puede borrar

% CONFIGURACIONES FINALES
% Template:     Informe/Reporte LaTeX
% Documento:    Configuraciones finales
% Versión:      4.6.2 (10/09/2017)
% Codificación: UTF-8
%
% Autor: Pablo Pizarro R.
%        Facultad de Ciencias Físicas y Matemáticas
%        Universidad de Chile
%        pablo.pizarro@ing.uchile.cl, ppizarror.com
%
% Manual template: [http://latex.ppizarror.com/Template-Informe/]
% Licencia MIT:    [https://opensource.org/licenses/MIT/]

% Se reestablecen headers y footers
\markboth{}{}
\newpage
\ifthenelse{\equal{\hfstyle}{style1}}{
	\fancyhead[L]{\nouppercase{\leftmark}}}{
}
\ifthenelse{\equal{\hfstyle}{style2}}{
	\fancyhead[L]{\nouppercase{\leftmark}}}{
}
\ifthenelse{\equal{\hfstyle}{style4}}{
	\fancyhead[L]{\nouppercase{\leftmark}}}{
}
\ifthenelse{\equal{\hfstyle}{style5}}{
	\fancyhead[R]{\nouppercase{\leftmark}}}{
}

% Se reestablece estilo de títulos
\sectionfont{\color{\titlecolor} \fontsizetitle \styletitle \selectfont}
\subsectionfont{\color{\subtitlecolor} \fontsizesubtitle \stylesubtitle \selectfont}
\subsubsectionfont{\color{\subsubtitlecolor} \fontsizesubsubtitle \stylesubsubtitle \selectfont}

% Se reestablecen números de página y secciones
\renewcommand{\thepage}{\arabic{page}}
\setcounter{page}{1}
\setcounter{section}{0}
\setcounter{footnote}{0}


% ========================= INICIO DEL DOCUMENTO =========================

% Template:     Informe/Reporte LaTeX
% Documento:    Archivo de ejemplo
% Versión:      6.1.1 (05/12/2018)
% Codificación: UTF-8
%
% Autor: Pablo Pizarro R. @ppizarror
%        Facultad de Ciencias Físicas y Matemáticas
%        Universidad de Chile
%        pablo.pizarro@ing.uchile.cl, ppizarror.com
%
% Manual template: [https://latex.ppizarror.com/Template-Informe/]
% Licencia MIT:    [https://opensource.org/licenses/MIT/]

% ------------------------------------------------------------------------------
% NUEVA SECCIÓN
% ------------------------------------------------------------------------------
% Las secciones se inician con \section, si se quiere una sección sin número se
% pueden usar las funciones \sectionanum (sección sin número) o la función
% \sectionanumnoi para crear el mismo título sin numerar y sin aparecer en el índice.
\section{Informes con \LaTeX}

	% SUB-SECCIÓN
	% Las sub-secciones se inician con \subsection, si se quiere una sub-sección sin
	% número se pueden usar las funciones \subsectionanum (nuevo subtítulo sin numeración)
	% o la función \subsectionanumnoi para crear el mismo subtítulo sin numerar y sin
	% aparecer en el índice
	\subsection{Una breve introducción}
		
		Este es un párrafo, puede contener múltiples \quotes{Expresiones} así como fórmulas o referencias\footnote{Las referencias se hacen utilizando la expresión \texttt{\textbackslash label}\{etiqueta\}.} a fórmulas como \eqref{eqn:identidad-imposible}. A continuación se muestra un ejemplo de inserción de imágenes o figuras (como la Figura \ref{img:testimage}) con el comando \href{https://latex.ppizarror.com/informe.html#hlp-imagen}{\texttt{\textbackslash insertimage}}:

		% Para insertar una imagen se puede usar la función \insertimage la cual
		% toma un primer parámetro opcional para definir una etiqueta (dentro de
		% los corchetes), luego toma la dirección de la imagen, sus parámetros
		% (en este caso se definió la escala de 0.15) y una leyenda opcional.
		\insertimage[\label{img:testimage}]{ejemplos/test-image.png}{scale=0.15}{Where are you? de \quotes{Internet}.}

		A continuación\footnote{Como se puede observar las funciones \texttt{\textbackslash insert...} añaden un párrafo automáticamente.} se muestra un ejemplo de inserción de ecuaciones simples con el comando \href{https://latex.ppizarror.com/informe.html#hlp-formulae}{\texttt{\textbackslash insertequation}}:

		% Se inserta una ecuación, el primer parámetro entre [] es opcional
		% (permite identificar con una etiqueta para poder referenciarlo después
		% con \ref), seguido de aquello se escribe la ecuación en modo bruto sin signos $.
		\insertequation[\label{eqn:identidad-imposible}]{\pow{a}{k}=\pow{b}{k}+\pow{c}{k} \quad \forall k>2}

		% Se añade párrafo de prueba. Notar que no se requiere añadir un salto
		% de línea después de insertar una ecuación.
		\lipsum[75]

		% Los párrafos se pueden añadir con \newp, esta función se hizo para
		% evitar errores y warnings por parte del compilador de LaTeX.
		\newp Este es un nuevo párrafo insertado con el comando \href{https://latex.ppizarror.com/informe.html#hlp-parrafo}{\texttt{\textbackslash newp}}. Si no te gustan los comandos \texttt{\textbackslash newp}, \texttt{\textbackslash newpar} o \texttt{\textbackslash newparnl} simplemente puedes usar los salto de línea convencionales acompañado de \texttt{\textbackslash par}.

	% SUB-SECCIÓN
	\subsection{Añadiendo tablas}

		También puedes usar tablas, ¡Crearlas es muy fácil!. Puedes usar el plugin \href{https://www.ctan.org/tex-archive/support/excel2latex/}{Excel2Latex} \cite{ref2} de Excel para convertir las tablas a \LaTeX\xspace o bien utilizar el \quotes{creador de tablas online} \cite{ref3}.

		% Tabla generada con el plugin Excel2Latex
		\begin{table}[htbp]
			\centering
			\caption{Ejemplo de tablas.}
			\begin{tabular}{ccc}
				\hline
				\textbf{Columna 1} & \textbf{Columna 2} & \textbf{Columna 3} \bigstrut\\
				\hline
				$\omega$ & $\nu$ & $\delta$ \bigstrut[t]\\
				$\beta$ & $\gamma$ & $\epsilon$ \\
				$\Phi$ & $\Theta$ & $\varSigma$ \bigstrut[b]\\
				\hline
			\end{tabular}
			\label{tab:tabla-1}
		\end{table}


% ------------------------------------------------------------------------------
% NUEVA SECCIÓN
% ------------------------------------------------------------------------------
\newpage
\section{Aquí un nuevo tema}

	% SUB-SECCIÓN
	\subsection{Haciendo informes como un profesional}

		% Se inserta una imagen flotante en la izquierda del documento con
		% \insertimageleft, al igual que las demás funciones, el primer parámetro
		% es opcional, luego viene la ubicación de la imagen, seguido de la escala
		% (un 30% del ancho de página) y por último su leyenda. Para insertar una
		% imagen flotante en la derecha se utiliza \insertimageright usando los
		% mismos parámetros.
		\insertimageleft[\label{img:imagen-izquierda}]{ejemplos/test-image-wrap}{0.3}{Apolo flotando a la izquierda.}

		\lipsum[1]

		% Párrafos de ejemplo
		\newp \lipsum[115]
		\newp \lipsum[2]

		% Agrega una ecuación con leyenda
		\insertequationcaptioned[\label{eqn:formulasinsentido}]{\int_{a}^{b} f(x) \dd{x} = \fracnpartial{f(x)}{x}{\eta} \cdotp \textstyle \sum_{x=a}^{b} f(x)\cancelto{1+\frac{\epsilon}{k}}{(1+\Delta x)}}{Ecuación sin sentido.}

		% Aquí no es necesario usar \newp dado que todas las funciones \insert... añaden un párrafo nuevo por defecto
		\lipsum[115]

		\newp \lipsum[4]

	% Inserta un subtítulo sin número
	\subsection{Otros párrafos más normales}

		% Párrafos
		\lipsum[7]
		\newp \lipsum[2]

		% Se inserta una ecuación larga con el entorno gathered (1 solo número de ecuación)
		\insertgathered[\label{eqn:eqn-larga}]{
			\lpow{\Lambda}{f} = \frac{L\cdot f}{W} \cdot \frac{\pow{\lpow{Q}{e}}{2}}{8 \pow{\pi}{2} \pow{W}{4} g} + \sum_{i=1}^{l} \frac{f \cdot \big( M - d\big)}{l \cdot W} \cdot \frac{\pow{\big(\lpow{Q}{e}- i\cdot Q\big)}{2}}{8 \pow{\pi}{2} \pow{W}{4} g}\\
			Q_e = 2.5Q \cdot \int_{0}^{e} V(x) \dd{x} + \aasin{ \bigg(1+\frac{1}{1-e}\bigg) }
		}

		% Nuevo párrafo
		\lipsum[4]

		% Se inserta un multicols, con esto se pueden escribir párrafos en varias columnas
		\begin{multicols}{2}

			% Párrafo 1
			\lipsum[4]

			% Ecuación encerrada en una caja
			\insertequation[]{ \boxed{f(x) = \fracdpartial{u}{t}} }

			% Párrafo 2 del multicols
			\lipsum[1]

		\end{multicols}

	% SUB-SECCIÓN
	\subsection{Ejemplos de inserción de código fuente}

		% A continuación se crea una función auxiliar, esta es una herramienta
		% extremadamente importante y muy útil. Esta función de ejemplo toma dos
		% parámetros, uno es el lenguaje del código fuente, el segundo el
		% identificador en el manual.
		\newcommand{\insertsrcmanual}[2]{\href{https://latex.ppizarror.com/informe.html\#hlp-srccode\&srctype=#1}{#2}}

		El template permite la inserción de los siguientes lenguajes de programación de forma nativa: \insertsrcmanual{bash}{bash}, \insertsrcmanual{c}{C}, \insertsrcmanual{csharp}{C\#}, \insertsrcmanual{cpp}{C++}, \insertsrcmanual{cuda}{cuda}, \insertsrcmanual{docker}{DOCKER}, \insertsrcmanual{html5}{HTML5}, \insertsrcmanual{java}{Java}, \insertsrcmanual{js}{Javascript}, \insertsrcmanual{json}{JSON}, \insertsrcmanual{kotlin}{Kotlin}, \insertsrcmanual{latex}{LaTeX}, \insertsrcmanual{matlab}{Matlab}, \insertsrcmanual{perl}{Perl}, \insertsrcmanual{php}{PHP}, \insertsrcmanual{plaintext}{Texto plano}, \insertsrcmanual{pseudocode}{Pseudocódigo}, \insertsrcmanual{python}{Python}, \insertsrcmanual{ruby}{Ruby}, \insertsrcmanual{scala}{Scala}, \insertsrcmanual{sql}{SQL} y \insertsrcmanual{xml}{XML}. Para insertar un código fuente se debe usar el entorno \texttt{sourcecode}, o el entorno \texttt{sourcecodep} si es que se quiere utilizar parámetros adicionales. \newp

		A continuación se presenta un ejemplo de inserción de código fuente en Python (Código \ref{codigo-python}), Java (Código \ref{codigo-java}) y Matlab (Código \ref{codigo-matlab}):

% Se define el lenguaje del código. Cuidado: Los códigos en LaTeX son sensibles a las tabulaciones y espacios en blanco
\begin{sourcecode}[\label{codigo-python}]{python}{Ejemplo en Python.}
import numpy as np

def incmatrix(genl1, genl2):
	m = len(genl1)
	n = len(genl2)
	M = None # Comentario 1
	VT = np.zeros((n*m, 1), int) # Comentario 2
\end{sourcecode}

\begin{sourcecode}[\label{codigo-java}]{java}{Ejemplo en Java.}
import java.io.IOException;
import javax.servlet.*;

// Hola mundo
public class Hola extends GenericServlet {
	public void service(ServletRequest request, ServletResponse response)
	throws ServletException, IOException{
		response.setContentType("text/html");
		PrintWriter pw = response.getWriter();
		pw.println("Hola, mundo!");
		pw.close();
	}
}
\end{sourcecode}

\begin{sourcecode}[\label{codigo-matlab}]{matlab}{Ejemplo en Matlab.}
% Se crea gráfico
f = figure(1); hold on;
movegui(f, 'center');
xlabel('td/Tn'); ylabel('FAD=Umax/Uf0');
title('Espectro de pulso de desplazamiento');

for j = 1:length(BETA)
	fad = ones(1, NDATOS); % Arreglo para el FAD
	for i = 1:NDATOS
		[t, u_t, ~, ~] = main(BETA(j), r(i), M, K, F0, 0);
		fad(i) = max(abs(u_t)) / uf0;
	end
end
\end{sourcecode}

	% SUB-SECCIÓN
	\subsection{Añadir múltiples imágenes}

	El template ofrece el entorno \href{https://latex.ppizarror.com/informe.html#hlp-images}{\texttt{images}} que permite insertar múltiples imágenes de una manera muy sencilla \footnote{Desde la versión \texttt{6.0.0} no se da soporte a las funciones de inserción de imágenes múltiples \texttt{\textbackslash insertdoubleimage}, \texttt{\textbackslash insertdoubleeqimage},\texttt{\textbackslash inserttripleimage}, \texttt{\textbackslash inserttripleeqimage}, \texttt{\textbackslash insertquadimage}, \texttt{\textbackslash insertpentaimage} y \texttt{\textbackslash inserthexaimage}.}. Para crear imágenes múltiples se deben usar las siguientes instrucciones:

\begin{sourcecode}{latex}{}
\begin{images}[\label{imagenmultiple}]{Ejemplo de imagen múltiple.}
	\addimage{ejemplos/test-image}{width=6.5cm}{Ciudad.}
	\addimage{ejemplos/test-image-wrap}{width=5cm}{Apolo.}
	\addimagenewline
	\addimage{ejemplos/test-image}{width=12cm}{Ciudad más grande.}
\end{images}
\end{sourcecode}

	Obteniendo así:

	\begin{images}{Ejemplo de imagen múltiple.}
		\addimage{ejemplos/test-image}{width=6.5cm}{Ciudad.}
		\addimage{ejemplos/test-image-wrap}{width=5cm}{Apolo.}
		\addimagenewline
		\addimage{ejemplos/test-image}{width=12cm}{Ciudad más grande.}
	\end{images}


% ------------------------------------------------------------------------------
% NUEVA SECCIÓN
% ------------------------------------------------------------------------------
% Inserta una sección sin número
\sectionanum{Más ejemplos}

	% Inserta un subtítulo sin número
	\subsectionanum{Listas y Enumeraciones}

		Hacer listas enumeradas con \LaTeX\ es muy fácil con el template\footnote{También puedes revisar el manual de las enumeraciones en \url{http://www.texnia.com/archive/enumitem.pdf}.}, para ello debes usar el comando \texttt{\textbackslash begin\{enumerate\}}, cada elemento comienza por \texttt{\textbackslash item}, resultando así:

		\begin{enumerate}
			\item Grecia
			\item Abracadabra
			\item Manzanas
		\end{enumerate}

		También se puede cambiar el tipo de enumeración, se pueden usar letras, números romanos, entre otros. Esto se logra cambiando el \textbf{label} del objeto \texttt{enumerate}. A continuación se muestra un ejemplo usando letras con el estilo \texttt{\textbackslash alph} \footnote{Con \texttt{\textbackslash Alph} las letras aparecen en mayúscula.}, números romanos con \texttt{\textbackslash roman} \footnote{Con \texttt{\textbackslash Roman} los números romanos salen en mayúscula.} o números griegos con \texttt{\textbackslash greek}\footnote{Una característica propia del template, con \texttt{\textbackslash Greek} las letras griegas están escritas en mayúscula.}:

		\begin{multicols}{3}
			\begin{enumeratebf}[label=\alph*) ] % Fuente en negrita
				\item Peras
				\item Manzanas
				\item Naranjas
			\end{enumeratebf}

			\begin{enumerate}[label=\greek*) ]
				\item Matemáticas
				\item Lenguaje
				\item Filosofía
			\end{enumerate}

			\begin{enumerate}[label=\roman*) ]
				\item Rojo
				\item Café
				\item Morado
			\end{enumerate}
		\end{multicols}

		Para hacer listas sin numerar con \LaTeX\ hay que usar el comando \texttt{\textbackslash begin\{itemize\}}, cada elemento empieza por \texttt{\textbackslash item}, resultando:

		\begin{multicols}{3}
			\begin{itemize}[label={--}]
				\item Peras
				\item Manzanas
				\item Naranjas
			\end{itemize}

			\begin{enumerate}[label={*}]
				\item Rojo
				\item Café
				\item Morado
			\end{enumerate}

			\begin{itemize}
				\item Árboles
				\item Pasto
				\item Flores
			\end{itemize}
		\end{multicols}

	% Inserta un subtítulo sin número
	\subsectionanum{Otros}

		Recuerda revisar el manual de todas las funciones y configuraciones de este template visitando el siguiente link: \url{https://latex.ppizarror.com/Template-Informe/}. Si necesitas una ayuda muy específica sobre el template, o si tienes alguna sugerencia, me puedes enviar un correo a \insertemail{pablo.pizarro@ing.uchile.cl}.


% ------------------------------------------------------------------------------
% REFERENCIAS (ESTILO BIBTEX), revisar configuración \stylecitereferences
% ------------------------------------------------------------------------------
\newpage % Salto de página
\begin{references}
	\bibitem{ref1}
	Template Informe en \LaTeX.
	\textit{¡Revisa el manual online de este template!} \\
	\url{https://latex.ppizarror.com/Template-Informe/}

	\bibitem{ref2}
	Excel2Latex.
	\textit{Importa de forma sencilla tus tablas de Excel a \LaTeX.} \\
	\url{https://www.ctan.org/tex-archive/support/excel2latex/}

	\bibitem{ref3}
	Overleaf.
	\textit{Uno de los mejores editores online para \LaTeX, renovado con su versión 2.0.} \\
	\href{https://es.overleaf.com}{\texttt{https://es.overleaf.com/}}
\end{references}


% ------------------------------------------------------------------------------
% ANEXO
% ------------------------------------------------------------------------------
\newpage
\begin{anexo}
	\section{Cálculos realizados}

		\subsection{Metodología}
			\lipsum[1]

			% Imagen, se numerará automáticamente con la letra del anexo según
			% la configuración \appendixindepobjnum
			\insertimage[\label{img:anexo-2}]{ejemplos/test-image.png}{scale=0.2}{Imagen en anexo.}

		\subsection{Resultados}
			\lipsum[10]

			% Tablas
			\begin{table}[htbp]
				\centering
				\caption{Tabla de cálculo.}
				\begin{tabular}{ccc}
					\hline
					\textbf{Elemento} & $\epsilon_i$ & \boldmath{}\textbf{Valor}\unboldmath{} \bigstrut\\
					\hline
					A     & 10    & 3,14$\pi$ \bigstrut[t]\\
					B     & 20    & 6 \\
					C     & 30    & 7 \\
					\end{tabular}
				\label{tab:anexo-1}
			\end{table}

	\newpage
	\section{Más cálculos}

		% Párrafo
		\lipsum[1]

		% Nuevo párrafo con identación
		\newp \lipsum[4]

		% Tabla de encuestas
		\begin{table}[htbp]
			\centering
			\caption{Resultados encuesta.}
			\begin{tabular}{ccc}
				\hline
				\textbf{Herramienta} & \textbf{Nota} & \textbf{Recomendado} \bigstrut\\
				\hline
				\LaTeX & 100\% & Si $\checkmark$ \\
				Microsoft Word \textsuperscript{\textregistered} & 0\% & No $\frownie$\\
			\end{tabular}
			\label{tab:anexo-2}
		\end{table}

\end{anexo}
 % Ejemplo, se puede borrar


% REFERENCIAS
\begin{references}
	\bibitem{ref1}
	Template Informe en \LaTeX.
	\textit{¡Revisa el manual online de este template!} \\
	\url{http://ppizarror.com/Template-Informe/}
	
	\bibitem{ref2}
	Excel2Latex.
	\textit{Importa de forma sencilla tus tablas de Excel a \LaTeX.} \\
	\url{https://www.ctan.org/tex-archive/support/excel2latex/}
	
	\bibitem{ref3}
	ShareLatex.
	\textit{Uno de los mejores editores online para \LaTeX\.} \\
	\href{https://www.sharelatex.com/?r=298b935f&rm=d&rs=b}{\texttt{http://www.tablesgenerator.com/}}
\end{references}

% FIN DEL DOCUMENTO
\end{document}